% Options for packages loaded elsewhere
% Options for packages loaded elsewhere
\PassOptionsToPackage{unicode}{hyperref}
\PassOptionsToPackage{hyphens}{url}
\PassOptionsToPackage{dvipsnames,svgnames,x11names}{xcolor}
%
\documentclass[
  letterpaper,
]{report}
\usepackage{xcolor}
\usepackage{amsmath,amssymb}
\setcounter{secnumdepth}{5}
\usepackage{iftex}
\ifPDFTeX
  \usepackage[T1]{fontenc}
  \usepackage[utf8]{inputenc}
  \usepackage{textcomp} % provide euro and other symbols
\else % if luatex or xetex
  \usepackage{unicode-math} % this also loads fontspec
  \defaultfontfeatures{Scale=MatchLowercase}
  \defaultfontfeatures[\rmfamily]{Ligatures=TeX,Scale=1}
\fi
\usepackage{lmodern}
\ifPDFTeX\else
  % xetex/luatex font selection
\fi
% Use upquote if available, for straight quotes in verbatim environments
\IfFileExists{upquote.sty}{\usepackage{upquote}}{}
\IfFileExists{microtype.sty}{% use microtype if available
  \usepackage[]{microtype}
  \UseMicrotypeSet[protrusion]{basicmath} % disable protrusion for tt fonts
}{}
\makeatletter
\@ifundefined{KOMAClassName}{% if non-KOMA class
  \IfFileExists{parskip.sty}{%
    \usepackage{parskip}
  }{% else
    \setlength{\parindent}{0pt}
    \setlength{\parskip}{6pt plus 2pt minus 1pt}}
}{% if KOMA class
  \KOMAoptions{parskip=half}}
\makeatother
% Make \paragraph and \subparagraph free-standing
\makeatletter
\ifx\paragraph\undefined\else
  \let\oldparagraph\paragraph
  \renewcommand{\paragraph}{
    \@ifstar
      \xxxParagraphStar
      \xxxParagraphNoStar
  }
  \newcommand{\xxxParagraphStar}[1]{\oldparagraph*{#1}\mbox{}}
  \newcommand{\xxxParagraphNoStar}[1]{\oldparagraph{#1}\mbox{}}
\fi
\ifx\subparagraph\undefined\else
  \let\oldsubparagraph\subparagraph
  \renewcommand{\subparagraph}{
    \@ifstar
      \xxxSubParagraphStar
      \xxxSubParagraphNoStar
  }
  \newcommand{\xxxSubParagraphStar}[1]{\oldsubparagraph*{#1}\mbox{}}
  \newcommand{\xxxSubParagraphNoStar}[1]{\oldsubparagraph{#1}\mbox{}}
\fi
\makeatother

\usepackage{color}
\usepackage{fancyvrb}
\newcommand{\VerbBar}{|}
\newcommand{\VERB}{\Verb[commandchars=\\\{\}]}
\DefineVerbatimEnvironment{Highlighting}{Verbatim}{commandchars=\\\{\}}
% Add ',fontsize=\small' for more characters per line
\usepackage{framed}
\definecolor{shadecolor}{RGB}{241,243,245}
\newenvironment{Shaded}{\begin{snugshade}}{\end{snugshade}}
\newcommand{\AlertTok}[1]{\textcolor[rgb]{0.68,0.00,0.00}{#1}}
\newcommand{\AnnotationTok}[1]{\textcolor[rgb]{0.37,0.37,0.37}{#1}}
\newcommand{\AttributeTok}[1]{\textcolor[rgb]{0.40,0.45,0.13}{#1}}
\newcommand{\BaseNTok}[1]{\textcolor[rgb]{0.68,0.00,0.00}{#1}}
\newcommand{\BuiltInTok}[1]{\textcolor[rgb]{0.00,0.23,0.31}{#1}}
\newcommand{\CharTok}[1]{\textcolor[rgb]{0.13,0.47,0.30}{#1}}
\newcommand{\CommentTok}[1]{\textcolor[rgb]{0.37,0.37,0.37}{#1}}
\newcommand{\CommentVarTok}[1]{\textcolor[rgb]{0.37,0.37,0.37}{\textit{#1}}}
\newcommand{\ConstantTok}[1]{\textcolor[rgb]{0.56,0.35,0.01}{#1}}
\newcommand{\ControlFlowTok}[1]{\textcolor[rgb]{0.00,0.23,0.31}{\textbf{#1}}}
\newcommand{\DataTypeTok}[1]{\textcolor[rgb]{0.68,0.00,0.00}{#1}}
\newcommand{\DecValTok}[1]{\textcolor[rgb]{0.68,0.00,0.00}{#1}}
\newcommand{\DocumentationTok}[1]{\textcolor[rgb]{0.37,0.37,0.37}{\textit{#1}}}
\newcommand{\ErrorTok}[1]{\textcolor[rgb]{0.68,0.00,0.00}{#1}}
\newcommand{\ExtensionTok}[1]{\textcolor[rgb]{0.00,0.23,0.31}{#1}}
\newcommand{\FloatTok}[1]{\textcolor[rgb]{0.68,0.00,0.00}{#1}}
\newcommand{\FunctionTok}[1]{\textcolor[rgb]{0.28,0.35,0.67}{#1}}
\newcommand{\ImportTok}[1]{\textcolor[rgb]{0.00,0.46,0.62}{#1}}
\newcommand{\InformationTok}[1]{\textcolor[rgb]{0.37,0.37,0.37}{#1}}
\newcommand{\KeywordTok}[1]{\textcolor[rgb]{0.00,0.23,0.31}{\textbf{#1}}}
\newcommand{\NormalTok}[1]{\textcolor[rgb]{0.00,0.23,0.31}{#1}}
\newcommand{\OperatorTok}[1]{\textcolor[rgb]{0.37,0.37,0.37}{#1}}
\newcommand{\OtherTok}[1]{\textcolor[rgb]{0.00,0.23,0.31}{#1}}
\newcommand{\PreprocessorTok}[1]{\textcolor[rgb]{0.68,0.00,0.00}{#1}}
\newcommand{\RegionMarkerTok}[1]{\textcolor[rgb]{0.00,0.23,0.31}{#1}}
\newcommand{\SpecialCharTok}[1]{\textcolor[rgb]{0.37,0.37,0.37}{#1}}
\newcommand{\SpecialStringTok}[1]{\textcolor[rgb]{0.13,0.47,0.30}{#1}}
\newcommand{\StringTok}[1]{\textcolor[rgb]{0.13,0.47,0.30}{#1}}
\newcommand{\VariableTok}[1]{\textcolor[rgb]{0.07,0.07,0.07}{#1}}
\newcommand{\VerbatimStringTok}[1]{\textcolor[rgb]{0.13,0.47,0.30}{#1}}
\newcommand{\WarningTok}[1]{\textcolor[rgb]{0.37,0.37,0.37}{\textit{#1}}}

\usepackage{longtable,booktabs,array}
\newcounter{none} % for unnumbered tables
\usepackage{calc} % for calculating minipage widths
% Correct order of tables after \paragraph or \subparagraph
\usepackage{etoolbox}
\makeatletter
\patchcmd\longtable{\par}{\if@noskipsec\mbox{}\fi\par}{}{}
\makeatother
% Allow footnotes in longtable head/foot
\IfFileExists{footnotehyper.sty}{\usepackage{footnotehyper}}{\usepackage{footnote}}
\makesavenoteenv{longtable}
\usepackage{graphicx}
\makeatletter
\newsavebox\pandoc@box
\newcommand*\pandocbounded[1]{% scales image to fit in text height/width
  \sbox\pandoc@box{#1}%
  \Gscale@div\@tempa{\textheight}{\dimexpr\ht\pandoc@box+\dp\pandoc@box\relax}%
  \Gscale@div\@tempb{\linewidth}{\wd\pandoc@box}%
  \ifdim\@tempb\p@<\@tempa\p@\let\@tempa\@tempb\fi% select the smaller of both
  \ifdim\@tempa\p@<\p@\scalebox{\@tempa}{\usebox\pandoc@box}%
  \else\usebox{\pandoc@box}%
  \fi%
}
% Set default figure placement to htbp
\def\fps@figure{htbp}
\makeatother
\usepackage{svg}





\setlength{\emergencystretch}{3em} % prevent overfull lines

\providecommand{\tightlist}{%
  \setlength{\itemsep}{0pt}\setlength{\parskip}{0pt}}



 


\makeatletter
\@ifpackageloaded{tcolorbox}{}{\usepackage[skins,breakable]{tcolorbox}}
\@ifpackageloaded{fontawesome5}{}{\usepackage{fontawesome5}}
\definecolor{quarto-callout-color}{HTML}{909090}
\definecolor{quarto-callout-note-color}{HTML}{0758E5}
\definecolor{quarto-callout-important-color}{HTML}{CC1914}
\definecolor{quarto-callout-warning-color}{HTML}{EB9113}
\definecolor{quarto-callout-tip-color}{HTML}{00A047}
\definecolor{quarto-callout-caution-color}{HTML}{FC5300}
\definecolor{quarto-callout-color-frame}{HTML}{acacac}
\definecolor{quarto-callout-note-color-frame}{HTML}{4582ec}
\definecolor{quarto-callout-important-color-frame}{HTML}{d9534f}
\definecolor{quarto-callout-warning-color-frame}{HTML}{f0ad4e}
\definecolor{quarto-callout-tip-color-frame}{HTML}{02b875}
\definecolor{quarto-callout-caution-color-frame}{HTML}{fd7e14}
\makeatother
\makeatletter
\@ifpackageloaded{bookmark}{}{\usepackage{bookmark}}
\makeatother
\makeatletter
\@ifpackageloaded{caption}{}{\usepackage{caption}}
\AtBeginDocument{%
\ifdefined\contentsname
  \renewcommand*\contentsname{Table of contents}
\else
  \newcommand\contentsname{Table of contents}
\fi
\ifdefined\listfigurename
  \renewcommand*\listfigurename{List of Figures}
\else
  \newcommand\listfigurename{List of Figures}
\fi
\ifdefined\listtablename
  \renewcommand*\listtablename{List of Tables}
\else
  \newcommand\listtablename{List of Tables}
\fi
\ifdefined\figurename
  \renewcommand*\figurename{Figure}
\else
  \newcommand\figurename{Figure}
\fi
\ifdefined\tablename
  \renewcommand*\tablename{Table}
\else
  \newcommand\tablename{Table}
\fi
}
\@ifpackageloaded{float}{}{\usepackage{float}}
\floatstyle{ruled}
\@ifundefined{c@chapter}{\newfloat{codelisting}{h}{lop}}{\newfloat{codelisting}{h}{lop}[chapter]}
\floatname{codelisting}{Listing}
\newcommand*\listoflistings{\listof{codelisting}{List of Listings}}
\makeatother
\makeatletter
\makeatother
\makeatletter
\@ifpackageloaded{caption}{}{\usepackage{caption}}
\@ifpackageloaded{subcaption}{}{\usepackage{subcaption}}
\makeatother
\usepackage{bookmark}
\IfFileExists{xurl.sty}{\usepackage{xurl}}{} % add URL line breaks if available
\urlstyle{same}
\hypersetup{
  pdftitle={Dark Web: De 0 a Experto},
  pdfauthor={Diego Saavedra},
  colorlinks=true,
  linkcolor={blue},
  filecolor={Maroon},
  citecolor={Blue},
  urlcolor={Blue},
  pdfcreator={LaTeX via pandoc}}


\title{Dark Web: De 0 a Experto}
\usepackage{etoolbox}
\makeatletter
\providecommand{\subtitle}[1]{% add subtitle to \maketitle
  \apptocmd{\@title}{\par {\large #1 \par}}{}{}
}
\makeatother
\subtitle{Análisis y estudio de la economía underground}
\author{Diego Saavedra}
\date{2025-12-31}
\begin{document}
\maketitle

\renewcommand*\contentsname{Table of contents}
{
\setcounter{tocdepth}{2}
\tableofcontents
}

\bookmarksetup{startatroot}

\chapter{Dark Web: De 0 a Experto}\label{dark-web-de-0-a-experto}

Análisis y estudio de la economía underground

\hfill\break

\bookmarksetup{startatroot}

\chapter{Dark Web: De 0 a Experto}\label{dark-web-de-0-a-experto-1}

\section{Análisis y estudio de la economía
underground}\label{anuxe1lisis-y-estudio-de-la-economuxeda-underground}

\begin{figure}[H]

{\centering \pandocbounded{\includesvg[keepaspectratio]{assets/images/logo-dark-web.svg}}

}

\caption{Logo del curso}

\end{figure}%

\begin{tcolorbox}[enhanced jigsaw, arc=.35mm, bottomrule=.15mm, bottomtitle=1mm, breakable, colback=white, colbacktitle=quarto-callout-note-color!10!white, colframe=quarto-callout-note-color-frame, coltitle=black, left=2mm, leftrule=.75mm, opacityback=0, opacitybacktitle=0.6, rightrule=.15mm, title=\textcolor{quarto-callout-note-color}{\faInfo}\hspace{0.5em}{Información del Curso}, titlerule=0mm, toprule=.15mm, toptitle=1mm]

\begin{itemize}
\tightlist
\item
  \textbf{Duración}: 12 semanas (3 créditos)
\item
  \textbf{Modalidad}: Presencial / Virtual sincrónico
\item
  \textbf{Requisitos previos}: Conocimientos básicos de redes, sistemas
  operativos y programación
\item
  \textbf{Herramientas}: Tor Browser, Whonix, entornos virtualizados
  aislados
\item
  \textbf{Advertencia}: Todo el material se basa en análisis de fuentes
  abiertas. El acceso a herramientas reales requiere entornos
  controlados y con autorización legal.
\end{itemize}

\end{tcolorbox}

\section{¿Qué aprenderás?}\label{quuxe9-aprenderuxe1s}

En este curso explorarás el ecosistema oculto de la dark web desde una
perspectiva técnica y de investigación, entendiendo cómo funcionan las
herramientas maliciosas, cómo se comercializan y cómo defenderse de
ellas.

\subsection{Competencias que
desarrollarás:}\label{competencias-que-desarrollaruxe1s}

\begin{enumerate}
\def\labelenumi{\arabic{enumi}.}
\tightlist
\item
  \textbf{Navegar de forma segura} por la dark web usando Tor y Whonix
\item
  \textbf{Analizar herramientas del underground} como loaders, stealers
  y ransomware
\item
  \textbf{Entender el modelo económico} de la cibercriminalidad
  organizada
\item
  \textbf{Aplicar técnicas de reverse engineering} a malware real
\item
  \textbf{Desarrollar estrategias de defensa} basadas en conocimiento de
  amenazas
\end{enumerate}

\section{¡Comencemos!}\label{comencemos}

\begin{quote}
``Conocer al enemigo es la mejor forma de prepararse para la
batalla.''\\
--- Sun Tzu, adaptado al contexto de ciberseguridad
\end{quote}

\bookmarksetup{startatroot}

\chapter{Sílabo - Dark Web: De 0 a
Experto}\label{suxedlabo---dark-web-de-0-a-experto}

\bookmarksetup{startatroot}

\chapter{Sílabo de la Asignatura}\label{suxedlabo-de-la-asignatura}

\section{1. Información General}\label{informaciuxf3n-general}

{\def\LTcaptype{none} % do not increment counter
\begin{longtable}[]{@{}ll@{}}
\toprule\noalign{}
Campo & Valor \\
\midrule\noalign{}
\endhead
\bottomrule\noalign{}
\endlastfoot
\textbf{Carrera} & Ingeniería en Sistemas/Computación \\
\textbf{Código} & DW101 \\
\textbf{Créditos} & 3 \\
\textbf{Horas Teoría} & 2 h/semana \\
\textbf{Horas Práctica} & 2 h/semana \\
\textbf{Horas Autónomo} & 3 h/semana \\
\textbf{Prerrequisitos} & Redes y Sistemas (RR001) \\
\textbf{Docente} & Diego Saavedra --- dsaavedra@unl.edu.ec \\
\textbf{Office Hours} & Miércoles 14:00-16:00 \\
\end{longtable}
}

\section{2. Descripción}\label{descripciuxf3n}

Este curso explora el ecosistema de la dark web desde una perseto
técnica de investigación en ciberseguridad. Los estudiantes analizarán
herramientas maliciosas, entenderán el modelo económico de la
cibercriminalidad y desarrollarán habilidades para defender
organizaciones contra amenazas emergentes. El enfoque combina teoría,
análisis técnico y laboratorios controlados.

\section{3. Resultados de Aprendizaje}\label{resultados-de-aprendizaje}

\subsection{Generales}\label{generales}

{\def\LTcaptype{none} % do not increment counter
\begin{longtable}[]{@{}
  >{\raggedright\arraybackslash}p{(\linewidth - 4\tabcolsep) * \real{0.2500}}
  >{\raggedright\arraybackslash}p{(\linewidth - 4\tabcolsep) * \real{0.3438}}
  >{\raggedright\arraybackslash}p{(\linewidth - 4\tabcolsep) * \real{0.4062}}@{}}
\toprule\noalign{}
\begin{minipage}[b]{\linewidth}\raggedright
Código
\end{minipage} & \begin{minipage}[b]{\linewidth}\raggedright
Resultado
\end{minipage} & \begin{minipage}[b]{\linewidth}\raggedright
Nivel Bloom
\end{minipage} \\
\midrule\noalign{}
\endhead
\bottomrule\noalign{}
\endlastfoot
OG1 & \textbf{Analizar} las técnicas de acceso a la dark web
\textbf{utilizando} Tor y Whonix, \textbf{verificando} la protección de
identidad & Analizar \\
OG2 & \textbf{Configurar} entornos de análisis de malware
\textbf{utilizando} VMs aisladads, \textbf{validando} la seguridad del
entorno & Aplicar \\
OG3 & \textbf{Evaluar} el impacto de la economía underground
\textbf{sobre} la seguridad organizacional, \textbf{justificando}
recomendaciones de defensa & Evaluar \\
OG4 & \textbf{Diseñar} estrategias de detección \textbf{basadas en}
técnicas de living-off-the-land, \textbf{demuestrando} eficacia &
Crear \\
\end{longtable}
}

\subsection{Específicos por Unidad}\label{especuxedficos-por-unidad}

{\def\LTcaptype{none} % do not increment counter
\begin{longtable}[]{@{}
  >{\raggedright\arraybackslash}p{(\linewidth - 6\tabcolsep) * \real{0.2000}}
  >{\raggedright\arraybackslash}p{(\linewidth - 6\tabcolsep) * \real{0.2000}}
  >{\raggedright\arraybackslash}p{(\linewidth - 6\tabcolsep) * \real{0.2750}}
  >{\raggedright\arraybackslash}p{(\linewidth - 6\tabcolsep) * \real{0.3250}}@{}}
\toprule\noalign{}
\begin{minipage}[b]{\linewidth}\raggedright
Unidad
\end{minipage} & \begin{minipage}[b]{\linewidth}\raggedright
Código
\end{minipage} & \begin{minipage}[b]{\linewidth}\raggedright
Resultado
\end{minipage} & \begin{minipage}[b]{\linewidth}\raggedright
Nivel Bloom
\end{minipage} \\
\midrule\noalign{}
\endhead
\bottomrule\noalign{}
\endlastfoot
U1 & OE1.1 & Listar las diferencias entre web clara, deep web y dark web
& Recordar \\
U1 & OE1.2 & Explicar el funcionamiento de la red Tor & Comprender \\
U2 & OE2.1 & Configurar un entorno Whonix para acceso seguro &
Aplicar \\
U3 & OE3.1 & Analizar un proof of concept de Log4Shell & Analizar \\
U4 & OE4.1 & Identificar técnicas de evasión en loaders & Recordar \\
U5 & OE5.1 & Examinar el código de un information stealer & Analizar \\
U6 & OE6.1 & Comparar banking trojans Android vs Windows & Analizar \\
U7 & OE7.1 & Desempacar malware usando técnicas estáticas & Aplicar \\
U8 & OE8.1 & Interpretar tráfico C2 de Cobalt Strike & Comprender \\
U9 & OE9.1 & Diagnosticar uso de LOLBins en un incidente & Analizar \\
U10 & OE10.1 & Evaluar la seguridad de un locker de ransomware &
Evaluar \\
U11 & OE11.1 & Comparar técnicas de ransomware Linux vs Windows &
Analizar \\
U12 & OE12.1 & Proponer mitigaciones para futuras amenazas & Crear \\
\end{longtable}
}

\section{4. Metodología}\label{metodologuxeda}

El curso sigue el enfoque ABACOM de \textbf{aprendizaje activo} con: -
\textbf{60\% práctica}: Laboratorios con ejecución manual de comandos -
\textbf{30\% teoría}: Conceptos, casos de estudio y análisis -
\textbf{10\% autónomo}: Lecturas, ejercicios y proyectos

Los laboratorios usan VMs completamente aisladadas. \textbf{No se
permite el copy-paste} --- cada comando se escribe manualmente para
construir competencia real.

\section{5. Evaluación}\label{evaluaciuxf3n}

\subsection{Matriz de Evaluación}\label{matriz-de-evaluaciuxf3n}

{\def\LTcaptype{none} % do not increment counter
\begin{longtable}[]{@{}
  >{\raggedright\arraybackslash}p{(\linewidth - 10\tabcolsep) * \real{0.1964}}
  >{\raggedright\arraybackslash}p{(\linewidth - 10\tabcolsep) * \real{0.1250}}
  >{\raggedright\arraybackslash}p{(\linewidth - 10\tabcolsep) * \real{0.2321}}
  >{\raggedright\arraybackslash}p{(\linewidth - 10\tabcolsep) * \real{0.1071}}
  >{\raggedright\arraybackslash}p{(\linewidth - 10\tabcolsep) * \real{0.1964}}
  >{\raggedright\arraybackslash}p{(\linewidth - 10\tabcolsep) * \real{0.1429}}@{}}
\toprule\noalign{}
\begin{minipage}[b]{\linewidth}\raggedright
Resultado
\end{minipage} & \begin{minipage}[b]{\linewidth}\raggedright
Pond.
\end{minipage} & \begin{minipage}[b]{\linewidth}\raggedright
Instrumento
\end{minipage} & \begin{minipage}[b]{\linewidth}\raggedright
Tipo
\end{minipage} & \begin{minipage}[b]{\linewidth}\raggedright
Criterios
\end{minipage} & \begin{minipage}[b]{\linewidth}\raggedright
Semana
\end{minipage} \\
\midrule\noalign{}
\endhead
\bottomrule\noalign{}
\endlastfoot
OE1.1-OE1.2 & 5\% & Lab 1: Introducción a Tor & Formativa & VM
funcionando, acceso exitoso & 2 \\
OE2.1 & 10\% & Quiz 2: OPSEC y Whonix & Sumativa & ≥ 70\% correctas &
2 \\
OE3.1 & 15\% & Lab 3: Análisis de exploits & Formativa & Identificación
correcta de vulnerabilidades & 4 \\
OE4.1 & 5\% & Quiz 4: Loaders & Formativa & ≥ 70\% correctas & 5 \\
OE5.1 & 5\% & Lab 5: Stealers & Formativa & Extracción correcta de IOCs
& 6 \\
OE6.1 & 5\% & Quiz 6: Banking trojans & Formativa & ≥ 70\% correctas &
7 \\
OE7.1 & 10\% & Lab 7: Unpacking & Formativa & Payload recuperado
exitosamente & 8 \\
OE8.1 & 10\% & Quiz 8: C2 Frameworks & Sumativa & ≥ 70\% correctas &
9 \\
OE9.1 & 5\% & Lab 9: LOLBins & Formativa & Técnicas detectadas
correctamente & 10 \\
OE10.1 & 10\% & Quiz 10: Ransomware & Sumativa & ≥ 70\% correctas &
11 \\
OE11.1 & 5\% & Lab 11: Linux ransomware & Formativa & Análisis
completado & 12 \\
OG1-OG4 & 20\% & Proyecto Final & Sumativa & Rúbrica adjunta & 13-15 \\
\end{longtable}
}

\section{6. Bibliografía}\label{bibliografuxeda}

\subsection{Básica}\label{buxe1sica}

\begin{enumerate}
\def\labelenumi{\arabic{enumi}.}
\tightlist
\item
  Kaye, L. (2026). \emph{Dissecting the Dark Web}. No Starch Press.
  ISBN: 978-1-7185-0461-5.
\item
  Saello, J. (2023). \emph{Practical Malware Analysis}. No Starch Press.
\item
  Kolias, V. et al.~(2022). \emph{Cybersecurity Education: A Hands-On
  Approach}. Springer.
\end{enumerate}

\subsection{Complementaria}\label{complementaria}

\begin{enumerate}
\def\labelenumi{\arabic{enumi}.}
\tightlist
\item
  MITRE ATT\&CK Framework --- https://attack.mitre.org
\item
  NIST SP 800-83 - Guide to Malware Incident Prevention and Handling
\item
  GitHub: DissectingTheDarkWeb -
  https://github.com/DissectingTheDarkWebBook
\end{enumerate}

\part{Unidad 1: Fundamentos de la Dark Web}

\chapter{Fundamentos de la Dark Web}\label{fundamentos-de-la-dark-web}

\chapter{Unidad 1: Fundamentos de la Dark
Web}\label{unidad-1-fundamentos-de-la-dark-web-1}

\section{1. Introducción}\label{introducciuxf3n}

¿Alguna vez te has preguntado cómo es posible que millones de datos
robados se comercien en la dark web? ¿Cómo los cibercriminales acceden a
estos mercados ocultos sin ser detectados? En 2021, el ataque a Colonial
Pipeline puso en evidencia la existencia de este ecosistema criminal.
Pero la dark web no es simplemente un ``lugar oscuro'' de internet; es
una red compleja con sus propias reglas, herramientas y economía.

Esta unidad te da la base para entender \textbf{qué es la dark web},
\textbf{cómo se diferencia de internet que conocemos}, y \textbf{por qué
es importante para la ciberseguridad}.

\subsection{Conexión al contexto}\label{conexiuxf3n-al-contexto}

La dark web representa el \textbf{3\% restante} del total de información
en internet. Mientras que internet visible (surface web) es accesible
mediante motores de búsqueda, la dark web requiere herramientas
especializadas y está compuesta principalmente por servicios criminales.

\section{2. Contenido Teórico}\label{contenido-teuxf3rico}

\subsection{2.1. La Web, el Deep Web y la Dark
Web}\label{la-web-el-deep-web-y-la-dark-web}

Los motores de búsqueda usan \textbf{crawlers} que indexan páginas. El
archivo \texttt{robots.txt} indica qué contenido indexar. La
\textbf{deep web} es el 95\%+ de internet no indexado (passwords, bases
deatos, paywalls). La \textbf{dark web} es el 3\% restante: sitios
\texttt{.onion} accesibles solo vía Tor.

\subsection{2.2. Onion Routing Detallado}\label{onion-routing-detallado}

El proceso usa \textbf{3 nodos} con capas de encriptación:

\begin{tcolorbox}[enhanced jigsaw, arc=.35mm, bottomrule=.15mm, bottomtitle=1mm, breakable, colback=white, colbacktitle=quarto-callout-tip-color!10!white, colframe=quarto-callout-tip-color-frame, coltitle=black, left=2mm, leftrule=.75mm, opacityback=0, opacitybacktitle=0.6, rightrule=.15mm, title=\textcolor{quarto-callout-tip-color}{\faLightbulb}\hspace{0.5em}{Tres nodos de Tor}, titlerule=0mm, toprule=.15mm, toptitle=1mm]

{\def\LTcaptype{none} % do not increment counter
\begin{longtable}[]{@{}lll@{}}
\toprule\noalign{}
Nodo & Conoce & No conoce \\
\midrule\noalign{}
\endhead
\bottomrule\noalign{}
\endlastfoot
\textbf{Entry Guard} & IP origen & Destino final \\
\textbf{Relay Node} & Nodo anterior/siguiente & IP origen o destino \\
\textbf{Exit Node} & Destino final & IP origen \\
\end{longtable}
}

\end{tcolorbox}

\textbf{Comandos de instalación Tor Browser (Linux)}:

\begin{Shaded}
\begin{Highlighting}[]
\FunctionTok{tar} \AttributeTok{{-}xf}\NormalTok{ tor\_browser\_archive.tar.xz}
\ExtensionTok{./start{-}tor{-}browser.desktop}
\end{Highlighting}
\end{Shaded}

\subsection{2.3. OPSEC esencial}\label{opsec-esencial}

\begin{itemize}
\tightlist
\item
  Usar Whonix-Gateway + Whonix-Workstation
\item
  VPN encima de Tor para ocultar uso de Tor
\item
  VM dedicada, nada de cuentas personales
\item
  Wipe/reimage fácil del sistema
\end{itemize}

\subsection{2.4. deepdarkCTI}\label{deepdarkcti}

Fuente para encontrar direcciones .onion: repositorio de inteligencia
abierta.

\section{3. Aplicaciones Prácticas}\label{aplicaciones-pruxe1cticas}

\subsection{Escenario 1: Verificación de
anonimato}\label{escenario-1-verificaciuxf3n-de-anonimato}

\begin{Shaded}
\begin{Highlighting}[]
\CommentTok{\# En Whonix}
\ExtensionTok{curl} \AttributeTok{{-}{-}socks5}\NormalTok{ 127.0.0.1:9050 https://check.torproject.org/api/ip}
\CommentTok{\# Output: \{"IsExit": true, "IP": "[IP EXCLUSIVA TOR]"\}}
\end{Highlighting}
\end{Shaded}

\subsection{Escenario 2: Búsqueda de
fuentes}\label{escenario-2-buxfasqueda-de-fuentes}

deepdarkCTI provee URLs .onion verificables para investigación.

\section{3. Aplicaciones Prácticas}\label{aplicaciones-pruxe1cticas-1}

\subsection{Escenario 1: Navegando foros
ocultos}\label{escenario-1-navegando-foros-ocultos}

Un investigador encuentra un foro .onion que vende credenciales robadas.
¿Cómo acceder sin exponer su identidad?

\textbf{Requisitos}: Whonix (gateway + workstation), encuesta de
amenazas, mecanismo de verificación.

\subsection{Escenario 2: Investigando precios de
exploits}\label{escenario-2-investigando-precios-de-exploits}

Un analista ve varios exploits vendidos por \$5,000 cada uno. ¿Qué
indica esto sobre el mercado?

\textbf{Análisis}: Los precios revelan la demanda, la complejidad y el
valor percibido.

\section{4. Caso de Estudio: El modelo Cyber Kill
Chain}\label{caso-de-estudio-el-modelo-cyber-kill-chain}

El libro utiliza el modelo \textbf{Cyber Kill Chain} de Lockheed Martin
para estructurar el ataque:

{\def\LTcaptype{none} % do not increment counter
\begin{longtable}[]{@{}
  >{\raggedright\arraybackslash}p{(\linewidth - 4\tabcolsep) * \real{0.1395}}
  >{\raggedright\arraybackslash}p{(\linewidth - 4\tabcolsep) * \real{0.3023}}
  >{\raggedright\arraybackslash}p{(\linewidth - 4\tabcolsep) * \real{0.5581}}@{}}
\toprule\noalign{}
\begin{minipage}[b]{\linewidth}\raggedright
Fase
\end{minipage} & \begin{minipage}[b]{\linewidth}\raggedright
Descripción
\end{minipage} & \begin{minipage}[b]{\linewidth}\raggedright
Herramientas en dark web
\end{minipage} \\
\midrule\noalign{}
\endhead
\bottomrule\noalign{}
\endlastfoot
1. Reconocimiento & Identificar objetivos & Herramientas de escaneo \\
2. Arma & Crear vector de ataque & Exploits, phishing kits \\
3. Entrega & Llevar arma al objetivo & Packers, loaders \\
4. Explotación & Ejecutar código malicioso & Exploit builders \\
5. Instalación & Persistir en el sistema & Backdoors, RATs \\
6. Comando y Control & Comunicación con víctima & C2 frameworks \\
7. Acción en Objetivo & Robo de datos, ransomware & Toolkits de
exfiltración \\
\end{longtable}
}

\section{5. Errores Comunes}\label{errores-comunes}

\subsection{❌ Misconception: ``Tor hace todo el trabajo de
anonimato''}\label{misconception-tor-hace-todo-el-trabajo-de-anonimato}

\textbf{Por qué es incorrecto}: Tor solo protege el tráfico de red. Si
envías tu nombre real, cookies de sesión previas, o usas el mismo
lenguaje de escritura, puedes ser identificado.

\textbf{✅ Comprensión correcta}: Tor es una herramienta que debes
combinar con buenas prácticas de OPSEC: VM aislada, nada de personal
identifiable, comportamiento neutro.

\subsection{❌ Misconception: ``La dark web es solo para
criminales''}\label{misconception-la-dark-web-es-solo-para-criminales}

\textbf{Por qué es incorrecto}: Periodistas, activistas, y personas bajo
regímenes opresivos usan la dark web para comunicación segura.

\textbf{✅ Comprensión correcta}: La dark web es una herramienta de
anonimato. Su uso depende de las intenciones del usuario.

\section{6. Actividades}\label{actividades}

\subsection{Lecturas Asignadas}\label{lecturas-asignadas}

\textbf{Obligatoria}: Kaye, L. (2026). \emph{Dissecting the Dark Web},
Capítulos 1-2, pp.~1-50.

\textbf{Opcional}:
\href{https://ssd.eff.org/en/module/how-tor-works}{``How Tor Works'' ---
Electronic Frontier Foundation}

\subsection{Ejercicios Prácticos}\label{ejercicios-pruxe1cticos}

\subsubsection{Ejercicio 1: Identificar tipos de
web}\label{ejercicio-1-identificar-tipos-de-web}

\textbf{Objetivo}: Diferenciar entre surface, deep y dark web.

\textbf{Pasos}: 1. Lista 5 ejemplos de surface web (ej: wikipedia.org)
2. Lista 5 ejemplos de deep web (ej: gmail.com) 3. Explica por qué cada
uno pertenece a su categoría

\textbf{Verifica}: Las diferencias clave están claramente identificadas.

\subsubsection{Ejercicio 2: Simular onion
routing}\label{ejercicio-2-simular-onion-routing}

\textbf{Objetivo}: Entender el concepto de enrutamiento en capas.

\textbf{Pasos}: 1. Escribe un mensaje de ejemplo 2. Encripta
``manualmente'' usando 3 claves diferentes (simulado) 3. Dibuena el
flujo de cómo cada nodo lo recibe y lo pasa

\textbf{Verifica}: El diagrama muestra claramente las capas y los nodos.

\subsection{Preguntas de Discusión}\label{preguntas-de-discusiuxf3n}

\begin{enumerate}
\def\labelenumi{\arabic{enumi}.}
\tightlist
\item
  ¿Deben los gobiernos prohibir tecnologías como Tor? ¿Por qué o por qué
  no?
\item
  ¿Cómo afecta el anonimato en la dark web al modelo de responsabilidad
  en ciberseguridad?
\end{enumerate}

\subsection{Auto-Evaluación}\label{auto-evaluaciuxf3n}

\subsection{Preguntas de
Verificación}\label{preguntas-de-verificaciuxf3n}

\begin{enumerate}
\def\labelenumi{\arabic{enumi}.}
\item
  \textbf{MCQ}: ¿Cuál capa del modelo OSI maneja la encriptación para
  Tor?

  \begin{enumerate}
  \def\labelenumii{\alph{enumii})}
  \tightlist
  \item
    Física
  \item
    Enlace
  \item
    Red
  \item
    Transporte
  \end{enumerate}

  \begin{quote}
  \textbf{Respuesta}: d) Transporte. Las conexiones TLS protegen la capa
  de transporte.
  \end{quote}
\item
  \textbf{Verdadero/Falso}: La dark web representa más del 50\% del
  contenido de internet.

  \begin{quote}
  \textbf{Falso}. La dark web es menos del 3\% del total.
  \end{quote}
\item
  \textbf{Corto}: ¿Qué significa ``onion routing'' en una frase?

  \begin{quote}
  \textbf{Respuesta}: Enrutamiento en capas donde cada nodo solo conoce
  su segmento de la ruta.
  \end{quote}
\end{enumerate}

\section{7. Bibliografía}\label{bibliografuxeda-1}

\subsection{Principal}\label{principal}

\begin{itemize}
\tightlist
\item
  Kaye, L. (2026). \emph{Dissecting the Dark Web}. No Starch Press.
  ISBN: 978-1-7185-0461-5.
\end{itemize}

\subsection{Papeles}\label{papeles}

\begin{itemize}
\tightlist
\item
  Anderson, R. (2021). ``The Tor Network: A Survey''. \emph{Journal of
  Cybersecurity}, 7(2). DOI: 10.1093/cybsec/tyab008
\end{itemize}

\subsection{Estándares}\label{estuxe1ndares}

\begin{itemize}
\tightlist
\item
  RFC 4949 --- Internet Security Glossary
\end{itemize}

\subsection{Digital}\label{digital}

\begin{itemize}
\tightlist
\item
  \href{https://tb-manual.torproject.org/}{Tor Project Documentation}
\item
  \href{https://attack.mitre.org/tactics/TA0001/}{MITRE ATT\&CK -
  Initial Access}
\end{itemize}

\section{8. Tareas Anti-IA}\label{tareas-anti-ia}

\subsection{Tarea 1: Dibujo manual}\label{tarea-1-dibujo-manual}

\textbf{Tarea}: En papel, dibuja el flujo de onion routing con 3 nodos.
Incluye: - Dirección IP de origen - Capas de encriptación - Qué ve cada
nodo

\textbf{Regla}: Sin herramientas digitales. Toma foto y sube.
Calificación: completitud, etiquetas, precisión.

\subsection{Tarea 2: Explicación
oral}\label{tarea-2-explicaciuxf3n-oral}

\textbf{Tarea}: Graba un video de 3 minutos explicando la diferencia
entre deep web y dark web como si enseñaras a un compañero.

\textbf{Requerimientos}: Usa analogías, responde 1 pregunta adicional
del instructor.

\textbf{Evaluación}: Claridad, corrección, capacidad de responder
preguntas.

\chapter{Laboratorio 1: Introducción a Tor y
Whonix}\label{laboratorio-1-introducciuxf3n-a-tor-y-whonix}

\chapter{Laboratorio 1: Introducción a Tor y
Whonix}\label{laboratorio-1-introducciuxf3n-a-tor-y-whonix-1}

\begin{tcolorbox}[enhanced jigsaw, arc=.35mm, bottomrule=.15mm, bottomtitle=1mm, breakable, colback=white, colbacktitle=quarto-callout-warning-color!10!white, colframe=quarto-callout-warning-color-frame, coltitle=black, left=2mm, leftrule=.75mm, opacityback=0, opacitybacktitle=0.6, rightrule=.15mm, title=\textcolor{quarto-callout-warning-color}{\faExclamationTriangle}\hspace{0.5em}{ADVERTENCIA DE SEGURIDAD}, titlerule=0mm, toprule=.15mm, toptitle=1mm]

Este laboratorio debe ejecutarse SOLO en: - Máquina virtual aislada
(Whonix, VirtualBox/VMware) - Red virtualhost-only - Sin acceso a redes
internas de la institución

\textbf{No ejecutar en máquina host}. Cualquier descarga o interacción
es bajo tu responsabilidad.

\end{tcolorbox}

\section{Objetivo}\label{objetivo}

Configurar y verificar un entorno Whonix seguro para exploración de la
dark web, entendiendo cada componente del sistema.

\section{Requisitos previos}\label{requisitos-previos}

\begin{itemize}
\tightlist
\item
  VirtualBox o VMware instalado
\item
  8GB RAM mínimo disponible
\item
  Conexión a internet estable
\end{itemize}

\section{Paso 1: Descargar Whonix}\label{paso-1-descargar-whonix}

\subsection{1.1. Identificar fuentes
oficiales}\label{identificar-fuentes-oficiales}

\textbf{Objetivo}: Encontrar la página oficial de Whonix.

\textbf{Pasos manuales}: 1. Abre una terminal en tu máquina host 2.
Ejecuta: \texttt{nslookup\ whonix.org} (NO copiar-pegar) 3. Verifica que
la IP resuelta coincida con documentación oficial

\textbf{Output esperado}:

\begin{verbatim}
Server:     8.8.8.8
Address:    8.8.8.8#53

Non-authoritative answer:
Name:   whonix.org
Address: [IP REAL]
\end{verbatim}

\subsection{1.1. DeepdarkCTI para URLs
.onion}\label{deepdarkcti-para-urls-.onion}

\textbf{Objetivo}: Identificar fuentes legítimas de URLs de dark web.

\textbf{Pasos manuales}: 1. Busca ``deepdarkCTI github'' en Google
(desde VM limpia) 2. Identifica el repositorio oficial 3. Lista URLs
.onion verificables en el README

\subsection{1.2. Verificar firma GPG}\label{verificar-firma-gpg}

\textbf{Objetivo}: Confirmar autenticidad de la ISO de Whonix.

\textbf{Pasos manuales}: 1. Descarga el archivo \texttt{.asc} de firma
2. Importa la clave: \texttt{gpg\ -\/-recv-keys\ {[}clave\ oficial{]}}
3. Verifica: \texttt{gpg\ -\/-verify\ {[}archivo.asc{]}}

\textbf{Output esperado}:

\begin{verbatim}
gpg: Good signature from "Whonix Release Signing Key"
\end{verbatim}

\section{Paso 2: Configurar Whonix
Gateway}\label{paso-2-configurar-whonix-gateway}

\subsection{2.1. Crear red host-only}\label{crear-red-host-only}

\textbf{Objetivo}: Aislar la VM de redes externas.

\textbf{Pasos manuales}:

\begin{Shaded}
\begin{Highlighting}[]
\CommentTok{\# En host (macOS/Linux)}
\ExtensionTok{VBoxManage}\NormalTok{ hostonlyif ipconfig vboxnet0 }\AttributeTok{{-}{-}ip}\NormalTok{ 192.168.56.1}
\end{Highlighting}
\end{Shaded}

\subsection{2.2. Configurar interfaz}\label{configurar-interfaz}

En VirtualBox: 1. Configura \textbf{Whonix-Gateway} con Adaptador 1: NAT
2. Configura \textbf{Whonix-Workstation} con Adaptador 1: Host-only

\section{Paso 3: Verificar funcionamiento de
Tor}\label{paso-3-verificar-funcionamiento-de-tor}

\subsection{3.1. Chequear servicio Tor}\label{chequear-servicio-tor}

\textbf{Objetivo}: Confirmar que Tor está corriendo.

\textbf{Pasos manuales}:

\begin{Shaded}
\begin{Highlighting}[]
\CommentTok{\# En Whonix Gateway}
\FunctionTok{sudo}\NormalTok{ systemctl status tor@default}
\end{Highlighting}
\end{Shaded}

\textbf{Output esperado}:

\begin{verbatim}
● tor@default.service - Anonymizing overlay network
   Loaded: loaded
   Active: active (running)
\end{verbatim}

\subsection{3.2. Verificar IP pública}\label{verificar-ip-puxfablica}

\textbf{Objetivo}: Confirmar anonimato.

\textbf{Pasos manuales}:

\begin{Shaded}
\begin{Highlighting}[]
\CommentTok{\# En Whonix Workstation}
\ExtensionTok{curl} \AttributeTok{{-}{-}socks5{-}hostname}\NormalTok{ 127.0.0.1:9050 https://check.torproject.org/api/ip}
\end{Highlighting}
\end{Shaded}

\textbf{Output esperado}:

\begin{Shaded}
\begin{Highlighting}[]
\FunctionTok{\{}\DataTypeTok{"IsExit"}\FunctionTok{:} \KeywordTok{true}\FunctionTok{,} \DataTypeTok{"IP"}\FunctionTok{:} \StringTok{"[IP DE SALIDA TOR]"}\FunctionTok{\}}
\end{Highlighting}
\end{Shaded}

\section{Paso 4: Navegador Tor
Browser}\label{paso-4-navegador-tor-browser}

\subsection{4.1. Configuración segura}\label{configuraciuxf3n-segura}

Verifica en Tor Browser: 1. Security Settings → \textbf{Safest} (máximo)
2. No instalar plugins 3. No guardar contraseñas 4. No descargar
archivos ejecutables

\subsection{4.2. Acceso a servicio de
prueba}\label{acceso-a-servicio-de-prueba}

\textbf{Objetivo}: Verificar navegación segura.

\textbf{Pasos manuales}: 1. Navega a
\texttt{http://check.torproject.org} 2. Verifica mensaje
``Congratulations. This browser is configured to use Tor.'' 3. Toma
captura de pantalla (solo la página de verificación)

\section{Verificación Final}\label{verificaciuxf3n-final}

{\def\LTcaptype{none} % do not increment counter
\begin{longtable}[]{@{}
  >{\raggedright\arraybackslash}p{(\linewidth - 4\tabcolsep) * \real{0.3000}}
  >{\raggedright\arraybackslash}p{(\linewidth - 4\tabcolsep) * \real{0.2250}}
  >{\raggedright\arraybackslash}p{(\linewidth - 4\tabcolsep) * \real{0.4750}}@{}}
\toprule\noalign{}
\begin{minipage}[b]{\linewidth}\raggedright
Checkpoint
\end{minipage} & \begin{minipage}[b]{\linewidth}\raggedright
Comando
\end{minipage} & \begin{minipage}[b]{\linewidth}\raggedright
Resultado esperado
\end{minipage} \\
\midrule\noalign{}
\endhead
\bottomrule\noalign{}
\endlastfoot
1 & \texttt{systemctl\ status\ tor@default} & active (running) \\
2 & \texttt{curl\ -\/-socks5\ 127.0.0.1:9050\ https://ipinfo.io/ip} & IP
diferente a tu real \\
3 & Tor Browser → check.torproject.org & Mensaje de confirmación \\
\end{longtable}
}

\section{Rubrica de laboratorio}\label{rubrica-de-laboratorio}

{\def\LTcaptype{none} % do not increment counter
\begin{longtable}[]{@{}
  >{\raggedright\arraybackslash}p{(\linewidth - 6\tabcolsep) * \real{0.1754}}
  >{\raggedright\arraybackslash}p{(\linewidth - 6\tabcolsep) * \real{0.2632}}
  >{\raggedright\arraybackslash}p{(\linewidth - 6\tabcolsep) * \real{0.1930}}
  >{\raggedright\arraybackslash}p{(\linewidth - 6\tabcolsep) * \real{0.3684}}@{}}
\toprule\noalign{}
\begin{minipage}[b]{\linewidth}\raggedright
Criterio
\end{minipage} & \begin{minipage}[b]{\linewidth}\raggedright
Excelente (5)
\end{minipage} & \begin{minipage}[b]{\linewidth}\raggedright
Bueno (3)
\end{minipage} & \begin{minipage}[b]{\linewidth}\raggedright
Necesita Mejora (1)
\end{minipage} \\
\midrule\noalign{}
\endhead
\bottomrule\noalign{}
\endlastfoot
Seguridad OPSEC & Todos los pasos seguros & 1-2 riesgos menores &
Multiple fallos de OPSEC \\
Ejecución manual & Todo ejecutado manualmente & Algunos copy-paste &
Mucha automatización \\
Documentación & Screenshots + notas claras & Screenshots básicos &
Incompleto \\
Verificación & Todas las IPs verificadas & Algunas verificadas & No
verificado \\
\end{longtable}
}

\chapter{Quiz 1: Fundamentos de la Dark
Web}\label{quiz-1-fundamentos-de-la-dark-web}

\chapter{Quiz 1: Fundamentos de la Dark
Web}\label{quiz-1-fundamentos-de-la-dark-web-1}

\section{Instrucciones}\label{instrucciones}

Responde las siguientes preguntas. Muestra tu trabajo donde sea
necesario. Usa tus propias palabras.

\begin{center}\rule{0.5\linewidth}{0.5pt}\end{center}

\section{Preguntas}\label{preguntas}

\subsection{1. Desarrollo (Bloom:
Comprender)}\label{desarrollo-bloom-comprender}

Explica en 3-5 oraciones la diferencia técnica entre el Surface Web y la
Dark Web. Incluye: - Cómo se indexa cada uno - Qué software se requiere
para acceder

\textbf{Puntos}: 10

\begin{center}\rule{0.5\linewidth}{0.5pt}\end{center}

\subsection{2. Opción múltiple (Bloom:
Recordar)}\label{opciuxf3n-muxfaltiple-bloom-recordar}

¿Cuántos nodos intermedios procesa el tráfico en la red Tor por defecto?

\begin{enumerate}
\def\labelenumi{\alph{enumi})}
\tightlist
\item
  1 nodo (entrada única)
\item
  3 nodos (entrada, medio, salida)
\item
  5 nodos (configuración máxima)
\item
  El número varia según carga
\end{enumerate}

\begin{quote}
\textbf{Respuesta}:
\end{quote}

\begin{center}\rule{0.5\linewidth}{0.5pt}\end{center}

\subsection{3. Verdadero/Falso (Bloom:
Evaluar)}\label{verdaderofalso-bloom-evaluar}

La deep web contiene exclusivamente contenido ilegal y está oculta por
malas intenciones.

\begin{quote}
\textbf{Respuesta}:\\
\textbf{Justificación}:
\end{quote}

\begin{center}\rule{0.5\linewidth}{0.5pt}\end{center}

\subsection{4. Análisis (Bloom:
Analizar)}\label{anuxe1lisis-bloom-analizar}

Dado el siguiente output de \texttt{curl} usando un proxy Tor:

\begin{verbatim}
{"IsExit": false, "IP": "185.220.101.1"}
\end{verbatim}

¿Qué puedes inferir? - ¿Estás usando Tor? - ¿Cuál es el IP de salida? -
¿Qué indica ``IsExit: false''?

\textbf{Puntos}: 10

\begin{center}\rule{0.5\linewidth}{0.5pt}\end{center}

\subsection{5. Comando (Bloom: Aplicar)}\label{comando-bloom-aplicar}

Escribe el comando para verificar tu IP usando el proxy SOCKS de Tor. No
copies el ejemplo del laboratorio --- escribe desde cero.

\textbf{Respuesta}:

\begin{center}\rule{0.5\linewidth}{0.5pt}\end{center}

\subsection{6. Diseño (Bloom: Crear)}\label{diseuxf1o-bloom-crear}

Dibuena (en papel) el flujo de información desde tu navegador hasta un
servidor .onion. Etiqueta: - Cliente - Entry guard - Relay nodes - Exit
node (si aplica) - Servidor destino

\textbf{Puntos}: 10

\begin{center}\rule{0.5\linewidth}{0.5pt}\end{center}

\section{Rubrica rápida}\label{rubrica-ruxe1pida}

{\def\LTcaptype{none} % do not increment counter
\begin{longtable}[]{@{}ll@{}}
\toprule\noalign{}
Pregunta & Puntos \\
\midrule\noalign{}
\endhead
\bottomrule\noalign{}
\endlastfoot
1 & 10 \\
2 & 5 \\
3 & 5 \\
4 & 10 \\
5 & 5 \\
6 & 10 \\
\textbf{Total} & \textbf{45} \\
\end{longtable}
}

Passing score: 32/45 (71\%)

\part{Unidad 2: Acceso Seguro y OPSEC}

\chapter{Acceso Seguro y OPSEC}\label{acceso-seguro-y-opsec}

\chapter{Unidad 2: Acceso Seguro y
OPSEC}\label{unidad-2-acceso-seguro-y-opsec-1}

\section{1. Introducción}\label{introducciuxf3n-1}

Acceder a la dark web sin el anonimato adecuado es como entrar a una
guerra sin equipo de protección. Las amenazas incluyen rastreo de IP,
malware, y hasta intervención legal. Esta unidad te enseña a protegerte.

\section{2. Contenido Teórico}\label{contenido-teuxf3rico-1}

\subsection{2.1. Principios de OPSEC}\label{principios-de-opsec}

Operational Security (OPSEC) es la disciplina de proteger información
sensible. Los principios clave:

\begin{enumerate}
\def\labelenumi{\arabic{enumi}.}
\tightlist
\item
  \textbf{Separación}: Ni tus datos personales ni tu comportamiento
  normal en la dark web
\item
  \textbf{Comportamiento neutral}: Escribe de forma diferente, actúa de
  forma diferente
\item
  \textbf{Infraestructura dedicada}: VMs limpias, no tu máquina
  principal
\end{enumerate}

\begin{tcolorbox}[enhanced jigsaw, arc=.35mm, bottomrule=.15mm, bottomtitle=1mm, breakable, colback=white, colbacktitle=quarto-callout-warning-color!10!white, colframe=quarto-callout-warning-color-frame, coltitle=black, left=2mm, leftrule=.75mm, opacityback=0, opacitybacktitle=0.6, rightrule=.15mm, title=\textcolor{quarto-callout-warning-color}{\faExclamationTriangle}\hspace{0.5em}{⚠️ Preguntas que nunca debes responder en dark web}, titlerule=0mm, toprule=.15mm, toptitle=1mm]

\begin{itemize}
\tightlist
\item
  ``¿Dónde vives?''
\item
  ``¿Cuál es tu trabajo real?''
\item
  ``¿Qué lenguaje escribo normalmente?''
\item
  ``¿Qué hora es en mi zona horaria?''
\end{itemize}

\end{tcolorbox}

\subsection{2.2. Arquitectura Whonix}\label{arquitectura-whonix}

Whonix tiene \textbf{dos VMs}:

{\def\LTcaptype{none} % do not increment counter
\begin{longtable}[]{@{}lll@{}}
\toprule\noalign{}
Componente & Función & Protección \\
\midrule\noalign{}
\endhead
\bottomrule\noalign{}
\endlastfoot
\textbf{Gateway} & Enrutamiento Tor & Aísla tráfico \\
\textbf{Workstation} & Navegador y herramientas & Sin IP directa \\
\end{longtable}
}

\begin{verbatim}
[Internet] → [Gateway Whonix] → [Workstation Whonix] → [Usuario]
\end{verbatim}

\subsection{2.3. Herramientas de defensa}\label{herramientas-de-defensa}

\begin{itemize}
\tightlist
\item
  \textbf{VM isolation}: VirtualBox con red host-only
\item
  \textbf{Network fingerprinting protection}: User agent randomizado
\item
  \textbf{Timing attacks}: No patrones predecibles
\item
  \textbf{Correlation resistance}: No doble uso de sesión
\end{itemize}

\section{3. Aplicaciones Prácticas}\label{aplicaciones-pruxe1cticas-2}

\subsection{Escenario: Investigador
identificado}\label{escenario-investigador-identificado}

Un investigador fue identificado porque: 1. Usó su frase típica en foros
2. Visitó sitios a la misma hora todos los días 3. Usó su navegador
principal con Tor

\subsection{Escenario: Setup seguro}\label{escenario-setup-seguro}

Un analista correcto: 1. VM completamente nueva 2. Comportamiento
impredecible 3. No personaliza nada

\section{4. Errores Comunes}\label{errores-comunes-1}

\subsection{❌ Misconception: ``Whonix es suficiente
solito''}\label{misconception-whonix-es-suficiente-solito}

\textbf{¿Por qué es falso?}: Whonix protege red, pero si usas mismo
idioma, horarios, o patrones, la correlación es posible.

\subsection{❌ Misconception: ``Puedo usar mi máquina para
esto''}\label{misconception-puedo-usar-mi-muxe1quina-para-esto}

\textbf{¿Por qué es falso?}: Cualquier instalación previa, cookies, o
configuraciones son huellas.

\section{5. Actividades}\label{actividades-1}

\subsection{Lectura}\label{lectura}

\begin{itemize}
\tightlist
\item
  Kaye, L. (2026). \emph{Dissecting the Dark Web}, Sección 1.2-1.5.
\end{itemize}

\subsection{Ejercicios}\label{ejercicios}

\subsubsection{Ejercicio 1: Matriz de
riesgos}\label{ejercicio-1-matriz-de-riesgos}

Lista 5 riesgos al acceder a dark web y propone mitigación.

\subsubsection{Ejercicio 2: Simulación de
timeline}\label{ejercicio-2-simulaciuxf3n-de-timeline}

Crea un horario impredecible para actividades de investigación.

\section{6. Auto-Evaluación}\label{auto-evaluaciuxf3n-1}

\begin{enumerate}
\def\labelenumi{\arabic{enumi}.}
\tightlist
\item
  \textbf{Definir}: ¿Qué significa operational security?
\item
  \textbf{Explicar}: ¿Por qué dos VMs en Whonix?
\item
  \textbf{Identificar}: Nombra 3 errores de OPSEC comunes.
\end{enumerate}

\section{7. Bibliografía}\label{bibliografuxeda-2}

\begin{itemize}
\tightlist
\item
  Kaye, L. (2026). \emph{Dissecting the Dark Web}, Capítulo 1.
\item
  NSA. (2020). \emph{OPSEC Guidelines}. DoD 5200.28.
\end{itemize}

\chapter{Laboratorio 2: Configuración Whonix
Avanzada}\label{laboratorio-2-configuraciuxf3n-whonix-avanzada}

\chapter{Laboratorio 2: Configuración Whonix
Avanzada}\label{laboratorio-2-configuraciuxf3n-whonix-avanzada-1}

\section{Objetivo}\label{objetivo-1}

Configurar Whonix con reglas de firewall adicionales y verificar la
protección contra correlation attacks.

\section{Paso 1: Configurar firewall
host}\label{paso-1-configurar-firewall-host}

\subsection{1.1. Reglas iptables}\label{reglas-iptables}

\begin{Shaded}
\begin{Highlighting}[]
\CommentTok{\# En host (macOS/Linux)}
\FunctionTok{sudo}\NormalTok{ iptables }\AttributeTok{{-}A}\NormalTok{ OUTPUT }\AttributeTok{{-}o}\NormalTok{ lo }\AttributeTok{{-}j}\NormalTok{ ACCEPT}
\FunctionTok{sudo}\NormalTok{ iptables }\AttributeTok{{-}A}\NormalTok{ OUTPUT }\AttributeTok{{-}m}\NormalTok{ conntrack }\AttributeTok{{-}{-}ctstate}\NormalTok{ ESTABLISHED }\AttributeTok{{-}j}\NormalTok{ ACCEPT}
\FunctionTok{sudo}\NormalTok{ iptables }\AttributeTok{{-}A}\NormalTok{ OUTPUT }\AttributeTok{{-}j}\NormalTok{ DROP}
\end{Highlighting}
\end{Shaded}

\section{Paso 2: Configurar Whonix
Gateway}\label{paso-2-configurar-whonix-gateway-1}

\subsection{2.1. Editar configuración
Tor}\label{editar-configuraciuxf3n-tor}

\begin{Shaded}
\begin{Highlighting}[]
\CommentTok{\# En Whonix Gateway}
\FunctionTok{sudo}\NormalTok{ nano /etc/tor/torrc}
\CommentTok{\# Agregar:}
\CommentTok{\# DisableNetworkTime=true}
\CommentTok{\# NewCircuitPeriod=600}
\end{Highlighting}
\end{Shaded}

\section{Paso 3: Verificar
aislamiento}\label{paso-3-verificar-aislamiento}

\subsection{3.1. Test de DNS leak}\label{test-de-dns-leak}

\begin{Shaded}
\begin{Highlighting}[]
\CommentTok{\# En Workstation}
\ExtensionTok{nslookup}\NormalTok{ google.com}
\CommentTok{\# Debe fallar o resolver vía Tor}
\end{Highlighting}
\end{Shaded}

\section{Verificación}\label{verificaciuxf3n}

{\def\LTcaptype{none} % do not increment counter
\begin{longtable}[]{@{}lll@{}}
\toprule\noalign{}
Checkpoint & Método & Resultado \\
\midrule\noalign{}
\endhead
\bottomrule\noalign{}
\endlastfoot
Firewall & \texttt{iptables\ -L} & Reglas aplicadas \\
DNS leak & \texttt{nslookup} & Sin leak \\
IP pública & Tor check & IP diferente \\
\end{longtable}
}

\section{Rubrica}\label{rubrica}

{\def\LTcaptype{none} % do not increment counter
\begin{longtable}[]{@{}llll@{}}
\toprule\noalign{}
Criterio & 5 pts & 3 pts & 1 pt \\
\midrule\noalign{}
\endhead
\bottomrule\noalign{}
\endlastfoot
Configuración & Todo correcto & 2 errores menores & Errores críticos \\
Verificación & 4/4 checkpoints & 2-3 checkpoints & \textless2
checkpoints \\
Documentación & Completa & Básica & Incompleta \\
\end{longtable}
}

\chapter{Quiz 2: OPSEC y Whonix}\label{quiz-2-opsec-y-whonix}

\chapter{Quiz 2: OPSEC y Whonix}\label{quiz-2-opsec-y-whonix-1}

\section{Preguntas}\label{preguntas-1}

\subsection{1. Desarrollo}\label{desarrollo}

¿Por qué es crítico usar una VM completamente nueva para navegar dark
web?

\subsection{2. Opción múltiple}\label{opciuxf3n-muxfaltiple}

¿Cuántas VMs tiene Whonix por defecto?

\begin{enumerate}
\def\labelenumi{\alph{enumi})}
\tightlist
\item
  1\\
\item
  2\\
\item
  3\\
\item
  Depende de la instalación
\end{enumerate}

\subsection{3. Verdadero/Falso}\label{verdaderofalso}

Los horarios predecibles de navegación pueden ser una señal de
identificación.

\subsection{4. Análisis}\label{anuxe1lisis}

Dado este log de firewall:

\begin{verbatim}
DROP tcp -- anywhere 192.168.1.100
ACCEPT tcp -- anywhere 10.152.152.10
\end{verbatim}

¿Qué política de red se está aplicando?

\subsection{5. Comando}\label{comando}

Escribe el comando para verificar el estado del servicio Tor.

\begin{center}\rule{0.5\linewidth}{0.5pt}\end{center}

\textbf{Puntuación total}: 30 pts \textbar{} \textbf{Passing}: 21 pts
(70\%)

\part{Unidad 3: Vulnerabilidades y Exploits}

\chapter{Vulnerabilidades, Exploits y Acceso
Inicial}\label{vulnerabilidades-exploits-y-acceso-inicial}

\chapter{Unidad 3: Vulnerabilidades, Exploits y Acceso
Inicial}\label{unidad-3-vulnerabilidades-exploits-y-acceso-inicial}

\section{1. Introducción}\label{introducciuxf3n-2}

La dark web funciona como un mercado donde el producto más valioso es el
\textbf{acceso}. Los Initial Access Brokers (IABs) venden puertas de
entrada a organizaciones, mientras que exploits como Log4Shell se
comercializan en paquetes listos para usar. Esta unidad te enseña a
entender estos productos desde una perspectiva defensiva.

\section{2. Contenido Teórico}\label{contenido-teuxf3rico-2}

\subsection{2.1. Initial Access Brokers
(IABs)}\label{initial-access-brokers-iabs}

Los IABs venden \textbf{acceso verificado} con precios basados en: -
Tamaño de la organización - Tipo de acceso (VPN, RDP, dominio) -
Información MFA incluida

{\def\LTcaptype{none} % do not increment counter
\begin{longtable}[]{@{}lll@{}}
\toprule\noalign{}
Tipo de acceso & Precio típico & Precio con MFA \\
\midrule\noalign{}
\endhead
\bottomrule\noalign{}
\endlastfoot
Personal & \$10-\$100 & \$50-\$200 \\
Corporativo pequeño & \$100-\$500 & \$300-\$800 \\
Credenciales dominio & \$500-\$2000 & \$1000-\$5000 \\
AWS/GCP/Azure & \$1000-\$5000 & \$2000-\$10000 \\
\end{longtable}
}

\subsection{2.2. Checkers vs
Brute-Forcers}\label{checkers-vs-brute-forcers}

{\def\LTcaptype{none} % do not increment counter
\begin{longtable}[]{@{}lll@{}}
\toprule\noalign{}
Herramienta & Función & Ejemplo del libro \\
\midrule\noalign{}
\endhead
\bottomrule\noalign{}
\endlastfoot
\textbf{Checkers} & Verifica credenciales existentes & TMChecker \\
\textbf{Brute-Forcers} & Fuerza bruta para nuevas & Forti VPN Bruter \\
\end{longtable}
}

Los checkers automatizan la verificación: - 15 millón de credenciales
sin verificar = \$400-\$700 - 7,000 credenciales verificadas
(corporativas + MFA) = \$1000+

\subsection{2.3. Exploit Builders}\label{exploit-builders}

No crean exploits nuevos --- son \textbf{mecanismos de entrega}: -
HTML-based phishing - LNK/URL attachments - Malvertising integrado

\textbf{Nota del libro}: Más baratos que exploits reales, requieren
menos skills técnicos.

\section{3. Caso de Estudio: TMChecker}\label{caso-de-estudio-tmchecker}

TMChecker verifica credenciales corporativas:

\begin{enumerate}
\def\labelenumi{\arabic{enumi}.}
\tightlist
\item
  \textbf{Entrada}: Lista de correos/usuarios
\item
  \textbf{Proceso}: Login automatizado a OWA/Office365
\item
  \textbf{Salida}: Archivo con credenciales válidas
\end{enumerate}

\subsection{Análisis técnico}\label{anuxe1lisis-tuxe9cnico}

\begin{Shaded}
\begin{Highlighting}[]
\CommentTok{\# Pseudocódigo del proceso de verificación}
\KeywordTok{def}\NormalTok{ verify\_credentials(username, password):}
\NormalTok{    session }\OperatorTok{=}\NormalTok{ requests.Session()}
\NormalTok{    response }\OperatorTok{=}\NormalTok{ session.post(}
        \StringTok{"https://outlook.office365.com/owa/"}\NormalTok{,}
\NormalTok{        data}\OperatorTok{=}\NormalTok{\{}\StringTok{"username"}\NormalTok{: username, }\StringTok{"password"}\NormalTok{: password\}}
\NormalTok{    )}
    \ControlFlowTok{if} \StringTok{"authenticated"} \KeywordTok{in}\NormalTok{ response.text:}
        \ControlFlowTok{return} \VariableTok{True}
    \ControlFlowTok{return} \VariableTok{False}
\end{Highlighting}
\end{Shaded}

\subsection{Señales de detección}\label{seuxf1ales-de-detecciuxf3n}

\begin{itemize}
\tightlist
\item
  Múltiples logins fallidos seguidos
\item
  User-Agent idéntico en todas las peticiones
\item
  Horarios predecibles de ataque
\end{itemize}

\section{4. Errores Comunes}\label{errores-comunes-2}

\subsection{❌ Misconception: ``Los exploits son imposibles de
usar''}\label{misconception-los-exploits-son-imposibles-de-usar}

\textbf{¿Por qué es falso?}: Los IABs empaquetan exploits con
instrucciones paso a paso. Baja la barrera de entrada.

\subsection{❌ Misconception: ``Phishing es
ineficaz''}\label{misconception-phishing-es-ineficaz}

\textbf{¿Por qué es falso?}: El 30\% de los ataques exitosos comienzan
con phishing. Es rentable.

\section{5. Actividades}\label{actividades-2}

\subsection{Lectura}\label{lectura-1}

\begin{itemize}
\tightlist
\item
  Kaye, L. (2026). \emph{Dissecting the Dark Web}, Capítulo 2.
\end{itemize}

\subsection{Ejercicios}\label{ejercicios-1}

\subsubsection{Ejercicio 1: Análisis de
CVE}\label{ejercicio-1-anuxe1lisis-de-cve}

Selecciona Log4Shell y explica: - Vector de ataque - Condiciones
necesarias - Mitigación

\subsubsection{Ejercicio 2: Identificar
IoCs}\label{ejercicio-2-identificar-iocs}

Dado un log de autenticación, identifica 3 indicadores de TMChecker.

\section{6. Auto-Evaluación}\label{auto-evaluaciuxf3n-2}

\begin{enumerate}
\def\labelenumi{\arabic{enumi}.}
\tightlist
\item
  \textbf{Listar}: 5 tipos de acceso vendidos por IABs
\item
  \textbf{Explicar}: Cómo verifica un checker las credenciales
\item
  \textbf{Analizar}: Diferencias entre Zphisher y Greatness kit
\end{enumerate}

\section{7. Bibliografía}\label{bibliografuxeda-3}

\begin{itemize}
\tightlist
\item
  MITRE ATT\&CK --- Initial Access Tactics
\item
  Kaye, L. (2026). \emph{Dissecting the Dark Web}, Capítulo 2.
\item
  NIST NVD --- CVE database.
\end{itemize}

\chapter{Laboratorio 3: Análisis de Credenciales y
Checkers}\label{laboratorio-3-anuxe1lisis-de-credenciales-y-checkers}

\chapter{Laboratorio 3: Análisis de Credenciales y
Checkers}\label{laboratorio-3-anuxe1lisis-de-credenciales-y-checkers-1}

\section{Objetivo}\label{objetivo-2}

Analizar el proceso de TMChecker para verificación de credenciales
corporativas.

\section{Paso 1: Identificar servicios
objetivo}\label{paso-1-identificar-servicios-objetivo}

\subsection{1.1. Servicios corporativos
comunes}\label{servicios-corporativos-comunes}

{\def\LTcaptype{none} % do not increment counter
\begin{longtable}[]{@{}lll@{}}
\toprule\noalign{}
Servicio & Puerto & Protocolo \\
\midrule\noalign{}
\endhead
\bottomrule\noalign{}
\endlastfoot
OWA/Office365 & 443 & HTTPS \\
VPN SSL & 443 & HTTPS \\
RDP & 3389 & TCP \\
SSH & 22 & TCP \\
\end{longtable}
}

\section{Paso 2: Identificar señales de
TMChecker}\label{paso-2-identificar-seuxf1ales-de-tmchecker}

\subsection{2.1. User-Agent estándar}\label{user-agent-estuxe1ndar}

\begin{Shaded}
\begin{Highlighting}[]
\CommentTok{\# TMChecker usa User{-}Agent predecible}
\ExtensionTok{curl} \AttributeTok{{-}H} \StringTok{"User{-}Agent: TMChecker/1.0"}\NormalTok{ target.com/login}
\end{Highlighting}
\end{Shaded}

\subsection{2.2. Patrones de login}\label{patrones-de-login}

\begin{Shaded}
\begin{Highlighting}[]
\CommentTok{\# Secuencial: user1, user2, user3...}
\CommentTok{\# O usando lista predefinida}
\CommentTok{\# Siempre mismo User{-}Agent}
\CommentTok{\# Siempre mismo timing entre requests}
\end{Highlighting}
\end{Shaded}

\section{Paso 3: Identificar en logs}\label{paso-3-identificar-en-logs}

\subsection{3.1. Buscar patrones}\label{buscar-patrones}

\begin{Shaded}
\begin{Highlighting}[]
\CommentTok{\# IPs múltiples, mismo User{-}Agent}
\FunctionTok{awk} \AttributeTok{{-}F}\StringTok{\textquotesingle{}"\textquotesingle{}} \StringTok{\textquotesingle{}/Login/\{print $6\}\textquotesingle{}}\NormalTok{ access.log }\KeywordTok{|} \FunctionTok{sort} \KeywordTok{|} \FunctionTok{uniq} \AttributeTok{{-}c}
\end{Highlighting}
\end{Shaded}

\section{Verificación}\label{verificaciuxf3n-1}

{\def\LTcaptype{none} % do not increment counter
\begin{longtable}[]{@{}ll@{}}
\toprule\noalign{}
Checkpoint & Resultado \\
\midrule\noalign{}
\endhead
\bottomrule\noalign{}
\endlastfoot
User-Agent & Identificado \\
Timing & Analizado \\
IP patterns & Detectado \\
\end{longtable}
}

\chapter{Quiz 3: Vulnerabilidades, Exploits y
Credenciales}\label{quiz-3-vulnerabilidades-exploits-y-credenciales}

\chapter{Quiz 3: Vulnerabilidades, Exploits y
Credenciales}\label{quiz-3-vulnerabilidades-exploits-y-credenciales-1}

\section{Preguntas}\label{preguntas-2}

\subsection{1. Desarrollo}\label{desarrollo-1}

Explica la diferencia entre un \textbf{checker} y un
\textbf{brute-forcer} con un ejemplo de cada uno.

\subsection{2. Opción múltiple}\label{opciuxf3n-muxfaltiple-1}

¿Cuál es el precio típico de credenciales corporativas VERIFICADAS (con
MFA) según el libro?

\begin{enumerate}
\def\labelenumi{\alph{enumi})}
\tightlist
\item
  \$100-\$500\\
\item
  \$500-\$1500\\
\item
  \$1000-\$5000\\
\item
  \$1500-\$3000
\end{enumerate}

\subsection{3. Verdadero/Falso}\label{verdaderofalso-1}

Los exploit builders crean vulnerabilidades nuevas desde cero.

\subsection{4. Análisis}\label{anuxe1lisis-1}

Si ves credenciales de 7,000 corporativas a \$1500+ y 15 millones de no
verificadas a \$700, ¿qué indica el precio diferencial?

\subsection{5. Comando}\label{comando-1}

Escribe el command para identificar User-Agent repetido en logs.

\subsection{6. Diseño}\label{diseuxf1o}

¿Qué diferencia hay entre un exploit real y un exploit builder? (máx 3
líneas)

\part{Unidad 4: Técnicas de Entrega de Malware}

\chapter{Técnicas de Entrega de
Malware}\label{tuxe9cnicas-de-entrega-de-malware}

\chapter{Unidad 4: Técnicas de Entrega de
Malware}\label{unidad-4-tuxe9cnicas-de-entrega-de-malware-1}

\section{1. Introducción}\label{introducciuxf3n-3}

Después de obtener acceso, el atacante necesita \textbf{entregar} el
payload. Los loaders son el primer estadio de ejecución; los botnets
proporcionan infraestructura masiva. Esta unidad analiza técnicas reales
de entrega.

\section{2. Loaders}\label{loaders}

\subsection{2.1. Qué es un loader}\label{quuxe9-es-un-loader}

Un loader \textbf{asegura la entrega} del payload. Pueden incluir: -
Keylogging - Screenshots - Ejecución de comandos - Funciones de
discovery

\subsection{2.2. SmokeLoader (desde 2013)}\label{smokeloader-desde-2013}

Técnicas de evasión identificadas: - \textbf{Opaque predicates}:
Expresiones que siempre son true/false pero difíciles de determinar -
\textbf{Function/string obfuscation}: Encripta funciones y strings para
ocultar propósito - \textbf{Anti-hooking}: Detecta y evita hooking de
sandboxes - \textbf{Anti-debugging}: Chequea flags de debugging
(\texttt{BeingDebugged}) - \textbf{Hash-based import resolution}: Oculta
funciones importadas usando hashes

\textbf{Precio en dark web}: \textasciitilde\$2,500 por builder que
genera HTA files

\subsection{2.3. BunnyLoader}\label{bunnyloader}

Características: - \textbf{File download y execution} -
\textbf{Anti-analysis técniques} (sleep, delays) - \textbf{Modular} con
plugins - \textbf{Registry persistence}

\section{3. Botnets}\label{botnets}

\subsection{3.1. Mirai Family}\label{mirai-family}

Mirai (2016) revolucionó el DDoS: - \textbf{Scanning}: Telnet en puertos
23/2323 - \textbf{Brute force}: Credenciales por defecto -
\textbf{Infection}: Binario MIPS/ARM

\subsection{3.2. CosaNostra Botnet}\label{cosanostra-botnet}

Botnet moderna con: - \textbf{C2 via Tor} - \textbf{DDoS + crypto
mining} - \textbf{Auto-update}

\section{4. Ejercicios}\label{ejercicios-2}

\begin{enumerate}
\def\labelenumi{\arabic{enumi}.}
\tightlist
\item
  Analiza un loader estático con strings y entropy
\item
  Identifica señales de Mirai en logs de red
\end{enumerate}

\section{5. Auto-Evaluación}\label{auto-evaluaciuxf3n-3}

\begin{enumerate}
\def\labelenumi{\arabic{enumi}.}
\tightlist
\item
  ¿Qué hace un loader?
\item
  Métodos anti-debug de SmokeLoader
\item
  Protocolo de infección de Mirai
\end{enumerate}

\chapter{Laboratorio 4: Análisis de
Loaders}\label{laboratorio-4-anuxe1lisis-de-loaders}

\chapter{Laboratorio 4: Análisis de
Loaders}\label{laboratorio-4-anuxe1lisis-de-loaders-1}

\section{Objetivo}\label{objetivo-3}

Identificar técnicas de evasión en SmokeLoader y BunnyLoader usando
análisis estático.

\section{Paso 1: SmokeLoader
Anti-debug}\label{paso-1-smokeloader-anti-debug}

\subsection{1.1. Chequeo de debugger}\label{chequeo-de-debugger}

\begin{Shaded}
\begin{Highlighting}[]
\NormalTok{; Assembly típico}
\NormalTok{call IsDebuggerPresent}
\NormalTok{test al, al}
\NormalTok{jz skip\_evasion}
\end{Highlighting}
\end{Shaded}

\section{Paso 2: BunnyLoader String
Obfuscation}\label{paso-2-bunnyloader-string-obfuscation}

\subsection{2.1. Identificar strings
encriptadas}\label{identificar-strings-encriptadas}

\begin{Shaded}
\begin{Highlighting}[]
\CommentTok{\# Entropy analysis}
\ExtensionTok{pecheck.py}\NormalTok{ malware.exe }\KeywordTok{|} \FunctionTok{grep} \AttributeTok{{-}A5} \AttributeTok{{-}B5} \StringTok{"entropy"}
\end{Highlighting}
\end{Shaded}

\section{Paso 3: Identificar rutas de
descarga}\label{paso-3-identificar-rutas-de-descarga}

\subsection{3.1. URLs ocultas}\label{urls-ocultas}

\begin{Shaded}
\begin{Highlighting}[]
\FunctionTok{strings}\NormalTok{ malware.exe }\KeywordTok{|} \FunctionTok{grep} \AttributeTok{{-}E} \StringTok{"https?://"}
\end{Highlighting}
\end{Shaded}

\section{Verificación}\label{verificaciuxf3n-2}

{\def\LTcaptype{none} % do not increment counter
\begin{longtable}[]{@{}ll@{}}
\toprule\noalign{}
Checkpoint & Resultado \\
\midrule\noalign{}
\endhead
\bottomrule\noalign{}
\endlastfoot
Anti-debug & Identificado \\
String enc & Rutas encontradas \\
URLs & Lista de C2 \\
\end{longtable}
}

\chapter{Quiz 4: Malware Delivery
Techniques}\label{quiz-4-malware-delivery-techniques}

\chapter{Quiz 4: Malware Delivery
Techniques}\label{quiz-4-malware-delivery-techniques-1}

\subsection{1. Desarrollo}\label{desarrollo-2}

¿Qué diferencia a un loader de un botnet?

\subsection{2. Opción múltiple}\label{opciuxf3n-muxfaltiple-2}

¿Qué técnica usa SmokeLoader para detectar depuración?

\begin{enumerate}
\def\labelenumi{\alph{enumi})}
\tightlist
\item
  Chequeo de IP pública\\
\item
  IsDebuggerPresent API\\
\item
  Análisis de tráfico\\
\item
  Fingerprint de navegador
\end{enumerate}

\subsection{3. Verdadero/Falso}\label{verdaderofalso-2}

Mirai se propaga mediante fuerza bruta de credenciales Telnet.

\subsection{4. Análisis}\label{anuxe1lisis-2}

¿Qué indica alta entropía en un binario?

\subsection{5. Comando}\label{comando-2}

Escribe el comando para identificar strings ofuscadas.

\part{Unidad 5: Information Stealers}

\chapter{Information Stealers}\label{information-stealers}

\chapter{Unidad 5: Information
Stealers}\label{unidad-5-information-stealers-1}

\section{1. Introducción}\label{introducciuxf3n-4}

Los stealers son el \textbf{columna vertebral} de la economía del
underground. Roban credenciales, cookies, wallets, y datos sensibles.
Este mercado genera miles de millones anuales.

\section{2. Proceso de un stealer}\label{proceso-de-un-stealer}

\begin{verbatim}
1. Harvesting → 2. Collection → 3. Exfiltration → 4. Monetización
\end{verbatim}

\subsection{2.1. Victim Download}\label{victim-download}

\begin{itemize}
\tightlist
\item
  Archivos ZIP
\item
  Documentos maliciosos
\item
  Cracks de software
\end{itemize}

\subsection{2.2. Data Collection}\label{data-collection}

{\def\LTcaptype{none} % do not increment counter
\begin{longtable}[]{@{}ll@{}}
\toprule\noalign{}
Tipo de dato & Ubicación típica \\
\midrule\noalign{}
\endhead
\bottomrule\noalign{}
\endlastfoot
Passwords navegadores & \texttt{AppData/Local/*/Login\ Data} \\
Cookies & \texttt{Cookies.sqlite}, \texttt{network.cookie} \\
Cryptowallets & \texttt{AppData/Roaming/*/wallet.dat} \\
Credit cards & Navegadores, keylogger \\
\end{longtable}
}

\subsection{2.3. Exfiltration}\label{exfiltration}

\begin{itemize}
\tightlist
\item
  \textbf{C2 domains}: \texttt{.ru}, \texttt{.tk}
\item
  \textbf{Telegram bots}: Envío directo
\item
  \textbf{FTP}: Servidores comprometidos
\end{itemize}

\section{3. Caso: Raccoon Stealer}\label{caso-raccoon-stealer}

\subsection{Funcionalidad clave}\label{funcionalidad-clave}

\begin{Shaded}
\begin{Highlighting}[]
\CommentTok{\# Pseudocódigo Raccoon}
\KeywordTok{def}\NormalTok{ collect\_browser\_data():}
\NormalTok{    browsers }\OperatorTok{=}\NormalTok{ [}\StringTok{"Chrome"}\NormalTok{, }\StringTok{"Firefox"}\NormalTok{, }\StringTok{"Edge"}\NormalTok{]}
    \ControlFlowTok{for}\NormalTok{ browser }\KeywordTok{in}\NormalTok{ browsers:}
        \CommentTok{\# Extraer logins}
\NormalTok{        logins }\OperatorTok{=}\NormalTok{ extract\_logins(browser)}
        \CommentTok{\# Extraer cookies}
\NormalTok{        cookies }\OperatorTok{=}\NormalTok{ extract\_cookies(browser)}
        \CommentTok{\# Extraer history}
\NormalTok{        history }\OperatorTok{=}\NormalTok{ extract\_history(browser)}
\end{Highlighting}
\end{Shaded}

\subsection{Evasion}\label{evasion}

\begin{itemize}
\tightlist
\item
  \textbf{Mutex}: Evita doble ejecución
\item
  \textbf{VM detection}: Salta en máquinas virtuales
\item
  \textbf{Sleep}: Delay para evitar sandbox
\end{itemize}

\section{4. Ejercicios}\label{ejercicios-3}

\begin{enumerate}
\def\labelenumi{\arabic{enumi}.}
\tightlist
\item
  Analizar strings de SapphireStealer
\item
  Identificar rutas de extracción de credenciales
\end{enumerate}

\section{5. Auto-Evaluación}\label{auto-evaluaciuxf3n-4}

\begin{enumerate}
\def\labelenumi{\arabic{enumi}.}
\tightlist
\item
  ¿Por qué siguen funcionando los stealers?
\item
  Métodos de exfiltración más comunes
\item
  Diferencias Raccoon vs RedLine
\end{enumerate}

\chapter{Laboratorio 5: Análisis de Information
Stealers}\label{laboratorio-5-anuxe1lisis-de-information-stealers}

\chapter{Laboratorio 5: Análisis de Information
Stealers}\label{laboratorio-5-anuxe1lisis-de-information-stealers-1}

\section{Objetivo}\label{objetivo-4}

Identificar rutas de extracción de credenciales en SapphireStealer
usando técnicas estáticas.

\section{Paso 1: Identificar rutas de
datos}\label{paso-1-identificar-rutas-de-datos}

\subsection{1.1. Rutas de navegadores}\label{rutas-de-navegadores}

\begin{Shaded}
\begin{Highlighting}[]
\FunctionTok{strings}\NormalTok{ SapphireStealer.exe }\KeywordTok{|} \FunctionTok{grep} \AttributeTok{{-}i} \AttributeTok{{-}E} \StringTok{"(appdata|chrome|firefox|edge|login)"}
\end{Highlighting}
\end{Shaded}

\section{Paso 2: Identificar rutas de
cryptowallets}\label{paso-2-identificar-rutas-de-cryptowallets}

\subsection{2.1. Wallets comunes}\label{wallets-comunes}

\begin{Shaded}
\begin{Highlighting}[]
\FunctionTok{strings}\NormalTok{ SapphireStealer.exe }\KeywordTok{|} \FunctionTok{grep} \AttributeTok{{-}i} \AttributeTok{{-}E} \StringTok{"(wallet|bitcoin|ethereum|crypto)"}
\end{Highlighting}
\end{Shaded}

\section{Paso 3: Identificar rutas de
exfiltración}\label{paso-3-identificar-rutas-de-exfiltraciuxf3n}

\subsection{3.1. URLs de C2}\label{urls-de-c2}

\begin{Shaded}
\begin{Highlighting}[]
\FunctionTok{strings}\NormalTok{ SapphireStealer.exe }\KeywordTok{|} \FunctionTok{grep} \AttributeTok{{-}i} \AttributeTok{{-}E} \StringTok{"(http|https|ftp|dns|telegram)"}
\end{Highlighting}
\end{Shaded}

\section{Verificación}\label{verificaciuxf3n-3}

{\def\LTcaptype{none} % do not increment counter
\begin{longtable}[]{@{}ll@{}}
\toprule\noalign{}
Checkpoint & Resultado esperado \\
\midrule\noalign{}
\endhead
\bottomrule\noalign{}
\endlastfoot
Browser paths & Wallets identificados \\
Wallet paths & Rutas extraídas \\
C2 URLs & Dominios encontrados \\
\end{longtable}
}

\chapter{Quiz 5: Information
Stealers}\label{quiz-5-information-stealers}

\chapter{Quiz 5: Information
Stealers}\label{quiz-5-information-stealers-1}

\subsection{1. Listar}\label{listar}

5 tipos de datos que roban los stealers.

\subsection{2. Explicar}\label{explicar}

¿Por qué siguen siendo efectivos los stealers?

\subsection{3. Analizar}\label{analizar}

¿Qué diferencia hay entre Raccoon y RedLine stealers?

\subsection{4. Identificar}\label{identificar}

Ruta típica de credenciales en Chrome.

\subsection{5. Comando}\label{comando-3}

Escribe el comando para buscar strings relacionadas a cryptowallets.

\part{Unidad 6: Banking Trojans}

\chapter{Banking Trojans}\label{banking-trojans}

\chapter{Unidad 6: Banking Trojans}\label{unidad-6-banking-trojans-1}

\section{1. Introducción}\label{introducciuxf3n-5}

Los banking trojans son \textbf{malware financiero} especializado en
robo de credenciales bancarias. Android usa Accessibility API para
overlay, Windows usa inyección de procesos.

\section{2. Android Banking Trojans}\label{android-banking-trojans}

\subsection{2.1. Accessibility API Abuse}\label{accessibility-api-abuse}

Los trojans solicitan permiso
\texttt{android.permission.BIND\_ACCESSIBILITY\_SERVICE} para: -
\textbf{Leer pantalla} - \textbf{Interceptar clicks} - \textbf{Injectar
overlays}

\subsection{2.2. Cerberus Overlay}\label{cerberus-overlay}

\begin{Shaded}
\begin{Highlighting}[]
\CommentTok{// Pseudocódigo de overlay malicioso}
\KeywordTok{public} \KeywordTok{class}\NormalTok{ OverlayService }\KeywordTok{extends}\NormalTok{ AccessibilityService }\OperatorTok{\{}
    \AttributeTok{@Override}
    \KeywordTok{public} \DataTypeTok{void} \FunctionTok{onAccessibilityEvent}\OperatorTok{(}\NormalTok{AccessibilityEvent event}\OperatorTok{)} \OperatorTok{\{}
        \BuiltInTok{String}\NormalTok{ packageName }\OperatorTok{=}\NormalTok{ event}\OperatorTok{.}\FunctionTok{getPackageName}\OperatorTok{().}\FunctionTok{toString}\OperatorTok{();}
        \ControlFlowTok{if} \OperatorTok{(}\NormalTok{packageName}\OperatorTok{.}\FunctionTok{equals}\OperatorTok{(}\StringTok{"com.bank.app"}\OperatorTok{))} \OperatorTok{\{}
            \FunctionTok{injectPhishingOverlay}\OperatorTok{();}
        \OperatorTok{\}}
    \OperatorTok{\}}
\OperatorTok{\}}
\end{Highlighting}
\end{Shaded}

\section{3. Windows Banking Trojans}\label{windows-banking-trojans}

\subsection{3.1. TrickBot Browser
Injection}\label{trickbot-browser-injection}

Técnicas: - \textbf{Process injection} en navegadores - \textbf{Function
hooking} API - \textbf{Web injects} para modificar páginas

\subsection{3.2. Cadenas de ataque}\label{cadenas-de-ataque}

\begin{verbatim}
1. Download → 2. Inject → 3. Hook → 4. Steal → 5. Exfiltrate
\end{verbatim}

\section{4. Mitigaciones}\label{mitigaciones}

\begin{itemize}
\tightlist
\item
  \textbf{Google Play Protect}
\item
  \textbf{EDR con detección de overlays}
\item
  \textbf{Políticas de no instalar APKs externos}
\end{itemize}

\section{5. Auto-Evaluación}\label{auto-evaluaciuxf3n-5}

\begin{enumerate}
\def\labelenumi{\arabic{enumi}.}
\tightlist
\item
  ¿Qué hace AccessibilityService?
\item
  Técnica de inyección de TrickBot
\item
  Señales de detección de overlays
\end{enumerate}

\chapter{Laboratorio 6: Banking Trojans
Android}\label{laboratorio-6-banking-trojans-android}

\chapter{Laboratorio 6: Banking Trojans
Android}\label{laboratorio-6-banking-trojans-android-1}

\section{Objetivo}\label{objetivo-5}

Analizar el uso de Accessibility API en trojans bancarios mediante
inspección estática.

\section{Paso 1: Identificar
permisos}\label{paso-1-identificar-permisos}

\subsection{1.1. Manifest analysis}\label{manifest-analysis}

\begin{Shaded}
\begin{Highlighting}[]
\ExtensionTok{apktool}\NormalTok{ d malware.apk}
\FunctionTok{grep} \AttributeTok{{-}A2} \AttributeTok{{-}B2} \StringTok{"android.permission.BIND\_ACCESSIBILITY\_SERVICE"}\NormalTok{ malware/AndroidManifest.xml}
\end{Highlighting}
\end{Shaded}

\section{Paso 2: Identificar
overlays}\label{paso-2-identificar-overlays}

\subsection{2.1. Overlay code}\label{overlay-code}

\begin{Shaded}
\begin{Highlighting}[]
\FunctionTok{grep} \AttributeTok{{-}r} \StringTok{"TYPE\_APPLICATION\_OVERLAY"}\NormalTok{ malware/smali/}
\FunctionTok{grep} \AttributeTok{{-}r} \StringTok{"WindowManager.LayoutParams"}\NormalTok{ malware/smali/}
\end{Highlighting}
\end{Shaded}

\section{Verificación}\label{verificaciuxf3n-4}

{\def\LTcaptype{none} % do not increment counter
\begin{longtable}[]{@{}ll@{}}
\toprule\noalign{}
Checkpoint & Resultado \\
\midrule\noalign{}
\endhead
\bottomrule\noalign{}
\endlastfoot
Accessibility perm & Identificado \\
Overlay code & Detectado \\
Suspicous perms & Listados \\
\end{longtable}
}

\chapter{Quiz 6: Banking Trojans}\label{quiz-6-banking-trojans}

\chapter{Quiz 6: Banking Trojans}\label{quiz-6-banking-trojans-1}

\subsection{1. Identificar}\label{identificar-1}

Permiso de Android que abusa Cerberus para overlays.

\subsection{2. Explicar}\label{explicar-1}

Técnica de process injection usada por TrickBot.

\subsection{3. Comparar}\label{comparar}

Banking trojan vs RAT principal diferencia.

\subsection{4. Verificar}\label{verificar}

Método para detectar overlays en Android.

\subsection{5. Comando}\label{comando-4}

APKtool command para decompilar un APK.

\part{Unidad 7: Packers y Crypters}

\chapter{Packers y Crypters}\label{packers-y-crypters}

\chapter{Unidad 7: Packers y
Crypters}\label{unidad-7-packers-y-crypters-1}

\section{1. Introducción}\label{introducciuxf3n-6}

Los packers \textbf{ofuscan} el malware para evadir AV/EDR. Los crypters
en el underground son herramientas de venta que demuestran
``funcionalidad''.

\section{2. Técnicas de protección}\label{tuxe9cnicas-de-protecciuxf3n}

\subsection{2.1. String Obfuscation}\label{string-obfuscation}

\begin{Shaded}
\begin{Highlighting}[]
\CommentTok{// String ofuscada}
\DataTypeTok{char} \OperatorTok{*}\NormalTok{enc }\OperatorTok{=} \StringTok{"}\SpecialCharTok{\textbackslash{}x48\textbackslash{}x65\textbackslash{}x6c\textbackslash{}x6c\textbackslash{}x6f}\StringTok{"}\OperatorTok{;} \CommentTok{// "Hello" XOR encrypted}
\ControlFlowTok{for}\OperatorTok{(}\DataTypeTok{int}\NormalTok{ i}\OperatorTok{=}\DecValTok{0}\OperatorTok{;}\NormalTok{ i}\OperatorTok{\textless{}}\NormalTok{len}\OperatorTok{;}\NormalTok{ i}\OperatorTok{++)}\NormalTok{ dec}\OperatorTok{[}\NormalTok{i}\OperatorTok{]} \OperatorTok{=}\NormalTok{ enc}\OperatorTok{[}\NormalTok{i}\OperatorTok{]} \OperatorTok{\^{}} \BaseNTok{0xAA}\OperatorTok{;}
\end{Highlighting}
\end{Shaded}

\subsection{2.2. Control Flow
Obfuscation}\label{control-flow-obfuscation}

\begin{itemize}
\tightlist
\item
  \textbf{Junk code insertion}
\item
  \textbf{Opaque predicates}
\item
  \textbf{Jump tables}
\end{itemize}

\subsection{2.3. Entropy Analysis}\label{entropy-analysis}

Entropy alta (\textgreater{} 7.0) indica: - Compresión - Encriptación\\
- Empaquetado

\section{3. Proceso de Unpacking}\label{proceso-de-unpacking}

\subsection{3.1. Memory Allocation}\label{memory-allocation}

\begin{enumerate}
\def\labelenumi{\arabic{enumi}.}
\tightlist
\item
  \texttt{VirtualAlloc} con permisos \texttt{PAGE\_EXECUTE\_READWRITE}
\item
  Copia del stub desempacador
\item
  Desencriptado del payload
\end{enumerate}

\subsection{3.2. UpX Unpacking}\label{upx-unpacking}

\begin{Shaded}
\begin{Highlighting}[]
\CommentTok{\# Identificar UpX}
\ExtensionTok{upx} \AttributeTok{{-}l}\NormalTok{ malware.exe}

\CommentTok{\# Desempaquetar}
\ExtensionTok{upx} \AttributeTok{{-}d}\NormalTok{ malware.exe }\AttributeTok{{-}o}\NormalTok{ unpacked.exe}
\end{Highlighting}
\end{Shaded}

\section{4. Ejercicios}\label{ejercicios-4}

\begin{enumerate}
\def\labelenumi{\arabic{enumi}.}
\tightlist
\item
  Calcular entropy de archivo binario
\item
  Identificar UPX con \texttt{strings} y \texttt{peid}
\end{enumerate}

\section{5. Auto-Evaluación}\label{auto-evaluaciuxf3n-6}

\begin{enumerate}
\def\labelenumi{\arabic{enumi}.}
\tightlist
\item
  Señales de archivo empaquetado
\item
  Método de unpacking manual
\item
  Uso de entropy en análisis
\end{enumerate}

\chapter{Laboratorio 7: Unpacking de
Malware}\label{laboratorio-7-unpacking-de-malware}

\chapter{Laboratorio 7: Unpacking de
Malware}\label{laboratorio-7-unpacking-de-malware-1}

\section{Objetivo}\label{objetivo-6}

Desempaquetar un binario protegido usando técnicas estáticas y
dinámicas.

\section{Paso 1: Identificar packer}\label{paso-1-identificar-packer}

\subsection{1.1. PEiD o Detect It Easy}\label{peid-o-detect-it-easy}

\begin{Shaded}
\begin{Highlighting}[]
\ExtensionTok{diec}\NormalTok{ malware.exe}
\CommentTok{\# O buscar firmas comunes}
\FunctionTok{strings}\NormalTok{ malware.exe }\KeywordTok{|} \FunctionTok{head} \AttributeTok{{-}20}
\end{Highlighting}
\end{Shaded}

\section{Paso 2: Calcular entropy}\label{paso-2-calcular-entropy}

\subsection{2.1. Análisis de entropía}\label{anuxe1lisis-de-entropuxeda}

\begin{Shaded}
\begin{Highlighting}[]
\CommentTok{\# Usar herramienta de entropía}
\ExtensionTok{entropyscan.py}\NormalTok{ malware.exe}
\CommentTok{\# Entropy alta (\textgreater{}7.0) indica compresión/encriptación}
\end{Highlighting}
\end{Shaded}

\section{Paso 3: Unpacking manual}\label{paso-3-unpacking-manual}

\subsection{3.1. Memory dump}\label{memory-dump}

En VMWare/VirtualBox: 1. Ejecutar en máquina aislada 2. Usar Process
Dump: \texttt{procdump\ -ma\ malware.exe} 3. Analizar dump con
IDA/Ghidra

\section{Verificación}\label{verificaciuxf3n-5}

{\def\LTcaptype{none} % do not increment counter
\begin{longtable}[]{@{}ll@{}}
\toprule\noalign{}
Checkpoint & Resultado \\
\midrule\noalign{}
\endhead
\bottomrule\noalign{}
\endlastfoot
Packer ID & Identificado \\
Entropy & Valor calculado \\
Payload & Recuperado \\
\end{longtable}
}

\chapter{Quiz 7: Packers y Crypters}\label{quiz-7-packers-y-crypters}

\chapter{Quiz 7: Packers y Crypters}\label{quiz-7-packers-y-crypters-1}

\subsection{1. Definir}\label{definir}

¿Qué hace un packer en malware?

\subsection{2. Identificar}\label{identificar-2}

Técnica para detectar malware empaquetado.

\subsection{3. Explicar}\label{explicar-2}

Método de unpacking estático vs dinámico.

\subsection{4. Diferenciar}\label{diferenciar}

Entropy alta significa qué en análisis.

\subsection{5. Comando}\label{comando-5}

Escribe el command para desempaquetar con UPX.

\part{Unidad 8: Frameworks C2}

\chapter{Frameworks
Command-and-Control}\label{frameworks-command-and-control}

\chapter{Unidad 8: Frameworks
Command-and-Control}\label{unidad-8-frameworks-command-and-control}

\section{1. Introducción}\label{introducciuxf3n-7}

Los frameworks C2 son \textbf{infraestructuras de control} para sistemas
comprometidos. Cobalt Strike es el estándar de facto, pero el
underground tiene variantes piratas.

\section{2. Componentes de un C2}\label{componentes-de-un-c2}

{\def\LTcaptype{none} % do not increment counter
\begin{longtable}[]{@{}ll@{}}
\toprule\noalign{}
Componente & Función \\
\midrule\noalign{}
\endhead
\bottomrule\noalign{}
\endlastfoot
\textbf{Implant} & Agente en víctima \\
\textbf{Beacon} & Comunicación peródica \\
\textbf{Team Server} & Control central \\
\textbf{Staging} & Configuración dinámica \\
\end{longtable}
}

\section{3. Caso: Vermilion Strike vs Cobalt
Strike}\label{caso-vermilion-strike-vs-cobalt-strike}

Vermilion Strike (copy de Cobalt Strike): - \textbf{Config decryption}:
Algoritmo XOR simple - \textbf{Metadata collection}: Sistema, usuario,
red - \textbf{Encrypted exfiltration}: AES-256

\subsection{Técnicas de evasión}\label{tuxe9cnicas-de-evasiuxf3n}

\begin{itemize}
\tightlist
\item
  \textbf{Sleep masks}: Cambian shellcode
\item
  \textbf{Process hollowing}: Inyección en proceso legítimo
\item
  \textbf{Named pipes}: Comunicación interna
\end{itemize}

\section{4. Detección}\label{detecciuxf3n}

\subsection{4.1. Tráfico C2}\label{truxe1fico-c2}

Señales en red: - Conexiones HTTP recurrentes - User-Agent sospechosos -
DNS tunneling

\subsection{4.2. Sliver Detection}\label{sliver-detection}

\begin{Shaded}
\begin{Highlighting}[]
\CommentTok{\# Análisis de tráfico}
\ExtensionTok{sliver{-}server} \AttributeTok{{-}s}  \CommentTok{\# Staging server logs}

\CommentTok{\# Detección de beacon}
\FunctionTok{netstat} \AttributeTok{{-}ano} \KeywordTok{|} \FunctionTok{grep}\NormalTok{ :443  }\CommentTok{\# Conexiones sospechosas}
\end{Highlighting}
\end{Shaded}

\section{5. Auto-Evaluación}\label{auto-evaluaciuxf3n-7}

\begin{enumerate}
\def\labelenumi{\arabic{enumi}.}
\tightlist
\item
  Arquitectura de Cobalt Strike
\item
  Configuración de beacon
\item
  Métodos de detección de C2
\end{enumerate}

\chapter{Laboratorio 8: Análisis de Frameworks
C2}\label{laboratorio-8-anuxe1lisis-de-frameworks-c2}

\chapter{Laboratorio 8: Análisis de Frameworks
C2}\label{laboratorio-8-anuxe1lisis-de-frameworks-c2-1}

\section{Objetivo}\label{objetivo-7}

Analizar configuración de Vermilion Strike/Cobalt Strike usando técnicas
de extracción.

\section{Paso 1: Identificar
configuración}\label{paso-1-identificar-configuraciuxf3n}

\subsection{1.1. Strings sospechosas}\label{strings-sospechosas}

\begin{Shaded}
\begin{Highlighting}[]
\FunctionTok{strings}\NormalTok{ malware.exe }\KeywordTok{|} \FunctionTok{grep} \AttributeTok{{-}i} \AttributeTok{{-}E} \StringTok{"(beacon|teamserver|callback|jarm)"}
\end{Highlighting}
\end{Shaded}

\section{Paso 2: Decodificar
configuración}\label{paso-2-decodificar-configuraciuxf3n}

\subsection{2.1. XOR decode}\label{xor-decode}

\begin{Shaded}
\begin{Highlighting}[]
\CommentTok{\# Pseudocódigo para decode}
\KeywordTok{def}\NormalTok{ decode\_config(data, key}\OperatorTok{=}\BaseNTok{0x55}\NormalTok{):}
    \ControlFlowTok{return} \BuiltInTok{bytes}\NormalTok{([b }\OperatorTok{\^{}}\NormalTok{ key }\ControlFlowTok{for}\NormalTok{ b }\KeywordTok{in}\NormalTok{ data])}
\end{Highlighting}
\end{Shaded}

\section{Verificación}\label{verificaciuxf3n-6}

{\def\LTcaptype{none} % do not increment counter
\begin{longtable}[]{@{}ll@{}}
\toprule\noalign{}
Checkpoint & Resultado \\
\midrule\noalign{}
\endhead
\bottomrule\noalign{}
\endlastfoot
Config strings & Extraídas \\
Decryption & Completada \\
C2 domains & Identificadas \\
\end{longtable}
}

\chapter{Quiz 8: Frameworks C2}\label{quiz-8-frameworks-c2}

\chapter{Quiz 8: Frameworks C2}\label{quiz-8-frameworks-c2-1}

\subsection{1. Identificar}\label{identificar-3}

Componentes principales de Cobalt Strike.

\subsection{2. Explicar}\label{explicar-3}

Diferencia entre Beacon y Implant.

\subsection{3. Analizar}\label{analizar-1}

Técnica de string deobfuscation en Vermilion Strike.

\subsection{4. Verificar}\label{verificar-1}

Método para detectar tráfico C2 en red.

\subsection{5. Comando}\label{comando-6}

Escribe el comando para listar conexiones de red activas.

\part{Unidad 9: Living Off the Land}

\chapter{Living Off the Land}\label{living-off-the-land}

\chapter{Unidad 9: Living Off the
Land}\label{unidad-9-living-off-the-land-1}

\section{1. Introducción}\label{introducciuxf3n-8}

``Living Off the Land'' significa usar \textbf{binarios legítimos} del
sistema para actividades maliciosas. PowerShell, WMI, y herramientas de
administrador se abusan para evadir detección.

\section{2. LOLBins comunes}\label{lolbins-comunes}

{\def\LTcaptype{none} % do not increment counter
\begin{longtable}[]{@{}lll@{}}
\toprule\noalign{}
Binario & Técnica típica & MITRE TID \\
\midrule\noalign{}
\endhead
\bottomrule\noalign{}
\endlastfoot
\texttt{powershell.exe} & Descarga y ejecución & T1059.001 \\
\texttt{wmic.exe} & Enumeración de procesos & T1047 \\
\texttt{rundll32.exe} & Ejecución de código & T1218.011 \\
\texttt{mshta.exe} & HTML applications maliciosas & T1218.005 \\
\texttt{wscript.exe} & Scripts JScript/VBScript & T1059.005 \\
\end{longtable}
}

\section{3. Técnicas por fase}\label{tuxe9cnicas-por-fase}

\subsection{3.1. Discovery}\label{discovery}

\begin{Shaded}
\begin{Highlighting}[]
\CommentTok{\# Enumeración de dominio}
\FunctionTok{Get{-}ADComputer} \OperatorTok{{-}}\NormalTok{Filter }\OperatorTok{*}
\NormalTok{arp }\OperatorTok{{-}}\NormalTok{a}
\NormalTok{net view}
\end{Highlighting}
\end{Shaded}

\subsection{3.2. Privilege Escalation}\label{privilege-escalation}

\begin{Shaded}
\begin{Highlighting}[]
\NormalTok{\# Abuso de servicios débiles}
\NormalTok{sc query }\KeywordTok{\textgreater{}}\NormalTok{ services.txt}
\NormalTok{icacls C:\textbackslash{}Windows\textbackslash{}Temp }\AttributeTok{/grant}\NormalTok{ Everyone:F}
\end{Highlighting}
\end{Shaded}

\subsection{3.3. Lateral Movement}\label{lateral-movement}

\begin{Shaded}
\begin{Highlighting}[]
\NormalTok{\# WMI para ejecución remota}
\NormalTok{wmic }\AttributeTok{/node:}\NormalTok{192.168.1.10 process call create }\StringTok{"malware.exe"}
\NormalTok{psexec \textbackslash{}\textbackslash{}192.168.1.10 }\AttributeTok{{-}u}\NormalTok{ admin }\AttributeTok{{-}p}\NormalTok{ passwd cmd.exe}
\end{Highlighting}
\end{Shaded}

\section{4. Caso: Conti Playbook}\label{caso-conti-playbook}

Documentos filtrados muestran: - Uso de \texttt{PsExec} para movimiento
lateral - \texttt{BloodHound} para mapeo de AD - \texttt{Rubeus} para
kerberoasting

\section{5. Mitigaciones}\label{mitigaciones-1}

\begin{itemize}
\tightlist
\item
  \textbf{Constrained Language Mode} para PowerShell
\item
  \textbf{AppLocker} para limitar binarios
\item
  \textbf{EDR con detección de comportamiento}
\item
  \textbf{Logging ampliado de procesos}
\end{itemize}

\section{6. Auto-Evaluación}\label{auto-evaluaciuxf3n-8}

\begin{enumerate}
\def\labelenumi{\arabic{enumi}.}
\tightlist
\item
  Lista 3 LOLBins y sus técnicas
\item
  Explica cómo funciona BloodHound
\item
  Propone mitigación para PsExec
\end{enumerate}

\chapter{Laboratorio 9: Living Off the
Land}\label{laboratorio-9-living-off-the-land}

\chapter{Laboratorio 9: Living Off the
Land}\label{laboratorio-9-living-off-the-land-1}

\section{Objetivo}\label{objetivo-8}

Identificar uso de LOLBins en logs de sistema mediante análisis de
eventos.

\section{Paso 1: Analizar logs de
proceso}\label{paso-1-analizar-logs-de-proceso}

\subsection{1.1. Sysmon events}\label{sysmon-events}

\begin{Shaded}
\begin{Highlighting}[]
\CommentTok{\textless{}!{-}{-} Evento 1: Process Create {-}{-}\textgreater{}}
\NormalTok{\textless{}}\KeywordTok{Event}\NormalTok{\textgreater{}}
\NormalTok{  \textless{}}\KeywordTok{EventData}\NormalTok{\textgreater{}}
\NormalTok{    \textless{}}\KeywordTok{Data} \OtherTok{Name=}\StringTok{"Image"}\NormalTok{\textgreater{}powershell.exe\textless{}/}\KeywordTok{Data}\NormalTok{\textgreater{}}
\NormalTok{    \textless{}}\KeywordTok{Data} \OtherTok{Name=}\StringTok{"CommandLine"}\NormalTok{\textgreater{}{-}enc [payload]\textless{}/}\KeywordTok{Data}\NormalTok{\textgreater{}}
\NormalTok{  \textless{}/}\KeywordTok{EventData}\NormalTok{\textgreater{}}
\NormalTok{\textless{}/}\KeywordTok{Event}\NormalTok{\textgreater{}}
\end{Highlighting}
\end{Shaded}

\section{Paso 2: Identificar patrones
sospechosos}\label{paso-2-identificar-patrones-sospechosos}

\subsection{2.1. PowerShell encoded}\label{powershell-encoded}

\begin{Shaded}
\begin{Highlighting}[]
\CommentTok{\# Buscar encoded commands}
\FunctionTok{grep} \AttributeTok{{-}i} \StringTok{"encodedcommand\textbackslash{}|enc "}\NormalTok{ sysmon.log}
\end{Highlighting}
\end{Shaded}

\section{Paso 3: WMI abuse detection}\label{paso-3-wmi-abuse-detection}

\subsection{3.1. WMIC patterns}\label{wmic-patterns}

\begin{Shaded}
\begin{Highlighting}[]
\FunctionTok{grep} \AttributeTok{{-}i} \StringTok{"wmic.*process.*call.*create"}\NormalTok{ sysmon.log}
\end{Highlighting}
\end{Shaded}

\section{Verificación}\label{verificaciuxf3n-7}

{\def\LTcaptype{none} % do not increment counter
\begin{longtable}[]{@{}ll@{}}
\toprule\noalign{}
Checkpoint & Detección \\
\midrule\noalign{}
\endhead
\bottomrule\noalign{}
\endlastfoot
Encoded PS & Identificado \\
WMI abuse & Detectado \\
LOLBin chain & Reconstruido \\
\end{longtable}
}

\chapter{Quiz 9: Living Off the Land}\label{quiz-9-living-off-the-land}

\chapter{Quiz 9: Living Off the
Land}\label{quiz-9-living-off-the-land-1}

\subsection{1. Listar}\label{listar-1}

5 LOLBins comunes en Windows.

\subsection{2. Explicar}\label{explicar-4}

Cómo BloodHound ayuda en lateral movement.

\subsection{3. Identificar}\label{identificar-4}

MITRE TID para PowerShell abuse.

\subsection{4. Verificar}\label{verificar-2}

Técnica defensiva contra PsExec.

\subsection{5. Comando}\label{comando-7}

Escribe el command para enumerar equipos en dominio.

\part{Unidad 10: Windows Ransomware}

\chapter{Windows Ransomware}\label{windows-ransomware}

\chapter{Unidad 10: Windows
Ransomware}\label{unidad-10-windows-ransomware-1}

\section{1. Introducción}\label{introducciuxf3n-9}

El ransomware es el \textbf{producto final} del pipeline de la dark web.
Lockers encriptan archivos y exigen rescate. El modelo
Ransomware-as-a-Service (RaaS) democratizó su uso.

\section{2. Historia del ransomware}\label{historia-del-ransomware}

{\def\LTcaptype{none} % do not increment counter
\begin{longtable}[]{@{}lll@{}}
\toprule\noalign{}
Etapa & Año & Característica \\
\midrule\noalign{}
\endhead
\bottomrule\noalign{}
\endlastfoot
Primeros & 2005-2012 & Encryption simples, envío por email \\
Modernos & 2013-2019 & Crypto, C2, evasión \\
RaaS & 2020+ & Afiliados, soporte, profit sharing \\
\end{longtable}
}

\section{3. Técnicas del locker}\label{tuxe9cnicas-del-locker}

\subsection{3.1. Generación de claves}\label{generaciuxf3n-de-claves}

\begin{Shaded}
\begin{Highlighting}[]
\CommentTok{// Algoritmo típico de generación}
\NormalTok{RSA }\OperatorTok{*}\NormalTok{rsa }\OperatorTok{=}\NormalTok{ RSA\_generate\_key}\OperatorTok{(}\DecValTok{2048}\OperatorTok{,}\NormalTok{ RSA\_F4}\OperatorTok{,}\NormalTok{ NULL}\OperatorTok{,}\NormalTok{ NULL}\OperatorTok{);}
\DataTypeTok{unsigned} \DataTypeTok{char}\NormalTok{ aes\_key}\OperatorTok{[}\DecValTok{32}\OperatorTok{];}
\NormalTok{RAND\_bytes}\OperatorTok{(}\NormalTok{aes\_key}\OperatorTok{,} \DecValTok{32}\OperatorTok{);}  \CommentTok{// AES{-}256 key for files}
\end{Highlighting}
\end{Shaded}

\subsection{3.2. Evasión de
recuperación}\label{evasiuxf3n-de-recuperaciuxf3n}

\begin{itemize}
\tightlist
\item
  \textbf{Delete shadow copies}: \texttt{vssadmin\ delete\ shadows}
\item
  \textbf{Exclude system dirs}: Evita que el OS no arranqué
\item
  \textbf{Whitelist archivos}: No encripta todo
\end{itemize}

\section{4. Caso: REvil/Sodinokibi}\label{caso-revilsodinokibi}

El libro compara \textbf{dos versiones}:

\subsection{Versión temprana (2019)}\label{versiuxf3n-temprana-2019}

\begin{itemize}
\tightlist
\item
  Encrypt simple con AES
\item
  Whitelist limitada: \texttt{C:\textbackslash{}Windows},
  \texttt{bootmgr}
\item
  Payload pequeño (\textasciitilde50KB)
\item
  Sin procesos adicionales
\end{itemize}

\subsection{Versión mejorada (2020+)}\label{versiuxf3n-mejorada-2020}

\begin{itemize}
\tightlist
\item
  \textbf{Process hollowing} para inyección
\item
  \textbf{Dynamic config} via comando de línea
\item
  \textbf{UAC bypass} mejorado
\item
  \textbf{Anti-analysis} integrado
\end{itemize}

\section{5. Snaffler para
reconocimiento}\label{snaffler-para-reconocimiento}

Antes de deployar ransomware, los atacantes buscan datos valiosos:

\begin{Shaded}
\begin{Highlighting}[]
\NormalTok{\# Enumerar shares accesibles}
\NormalTok{net view}
\NormalTok{net view \textbackslash{}\textbackslash{}target{-}system}
\NormalTok{net share}

\NormalTok{\# Snaffler para triage}
\NormalTok{Snaffler.exe }\AttributeTok{{-}s} \AttributeTok{{-}a} \AttributeTok{{-}o}\NormalTok{ shares.log}
\NormalTok{\# Salida: \{Green\} (bajo priority}\KeywordTok{)}\NormalTok{ vs \{Black\} (VM, SQL, passwords}\KeywordTok{)}
\end{Highlighting}
\end{Shaded}

\textbf{Nota de seguridad}: Snaffler es detectado por Defender como no
``OPSEC safe''.

\section{6. Decryptores}\label{decryptores}

Crean cuando: - Claves privadas filtradas (REvil shutdown 2021) -
Implementación criptográfica débil - Bug en RNG o key derivation

\subsection{Ejemplo: Babuk leak}\label{ejemplo-babuk-leak}

El source code filtrado permitió crear decryptor después de encontrar
errores.

\section{5. Decryptores}\label{decryptores-1}

Desarrolladores crean decryptores cuando: - Claves privadas filtradas -
Implementación débil - Colisiones de RNG

\section{6. Auto-Evaluación}\label{auto-evaluaciuxf3n-9}

\begin{enumerate}
\def\labelenumi{\arabic{enumi}.}
\tightlist
\item
  Fases del pipeline de ransomware
\item
  Técnicas de eliminación de backups
\item
  Cómo funciona el modelo RaaS
\end{enumerate}

\chapter{Laboratorio 10: Análisis de Ransomware
Windows}\label{laboratorio-10-anuxe1lisis-de-ransomware-windows}

\chapter{Laboratorio 10: Análisis de Ransomware
Windows}\label{laboratorio-10-anuxe1lisis-de-ransomware-windows-1}

\section{Objetivo}\label{objetivo-9}

Analizar técnicas de encriptación en REvil/Sodinokibi sin ejecutar
código.

\section{Paso 1: Identificar rutas
excluidas}\label{paso-1-identificar-rutas-excluidas}

\subsection{1.1. Whitelist strings}\label{whitelist-strings}

\begin{Shaded}
\begin{Highlighting}[]
\FunctionTok{strings}\NormalTok{ ransomware.exe }\KeywordTok{|} \FunctionTok{grep} \AttributeTok{{-}i} \AttributeTok{{-}E} \StringTok{"(window|system32|boot|whitelist)"}
\end{Highlighting}
\end{Shaded}

\section{Paso 2: Identificar rutas de eliminación de
backups}\label{paso-2-identificar-rutas-de-eliminaciuxf3n-de-backups}

\subsection{2.1. Shadow copy deletion}\label{shadow-copy-deletion}

\begin{Shaded}
\begin{Highlighting}[]
\FunctionTok{strings}\NormalTok{ ransomware.exe }\KeywordTok{|} \FunctionTok{grep} \AttributeTok{{-}i} \AttributeTok{{-}E} \StringTok{"(vssadmin|shadow|delete|wmic)"}
\end{Highlighting}
\end{Shaded}

\section{Paso 3: Identificar rutas de
exfiltración}\label{paso-3-identificar-rutas-de-exfiltraciuxf3n-1}

\subsection{3.1. Note drop}\label{note-drop}

\begin{Shaded}
\begin{Highlighting}[]
\FunctionTok{strings}\NormalTok{ ransomware.exe }\KeywordTok{|} \FunctionTok{grep} \AttributeTok{{-}i} \AttributeTok{{-}E} \StringTok{"(readme|decrypt|tor.*onion|payment)"}
\end{Highlighting}
\end{Shaded}

\section{Verificación}\label{verificaciuxf3n-8}

{\def\LTcaptype{none} % do not increment counter
\begin{longtable}[]{@{}ll@{}}
\toprule\noalign{}
Checkpoint & Resultado \\
\midrule\noalign{}
\endhead
\bottomrule\noalign{}
\endlastfoot
Whitelist & Identificado \\
VSS deletion & Detectado \\
Note drop & Encontrado \\
\end{longtable}
}

\chapter{Quiz 10: Windows Ransomware}\label{quiz-10-windows-ransomware}

\chapter{Quiz 10: Windows
Ransomware}\label{quiz-10-windows-ransomware-1}

\subsection{1. Identificar}\label{identificar-5}

Algoritmo típico para generar claves de ransomware.

\subsection{2. Explicar}\label{explicar-5}

Qué hace \texttt{vssadmin\ delete\ shadows} en ataque.

\subsection{3. Analizar}\label{analizar-2}

Diferencia entre versión temprana y mejorada de REvil.

\subsection{4. Verificar}\label{verificar-3}

Método para identificar ransomware en logs.

\subsection{5. Comando}\label{comando-8}

Escribe el command para matar procesos específicos.

\part{Unidad 11: Linux y ESXi Ransomware}

\chapter{Linux y ESXi Ransomware}\label{linux-y-esxi-ransomware}

\chapter{Unidad 11: Linux y ESXi
Ransomware}\label{unidad-11-linux-y-esxi-ransomware-1}

\section{1. Introducción}\label{introducciuxf3n-10}

Linux y ESXi son \textbf{blancos atractivos} porque: - Menos EDR que
Windows - Acceso a VMs completas - Data centers completos comprometidos

\section{2. ESXi Hypervisor}\label{esxi-hypervisor}

\subsection{2.1. Ventajas para atacantes}\label{ventajas-para-atacantes}

\begin{itemize}
\tightlist
\item
  VMs completas = más datos
\item
  Sin antivirus tradicional
\item
  Admin tool exposure (ports 80/443/902)
\end{itemize}

\subsection{2.2. REvil ESXi variant}\label{revil-esxi-variant}

\subsection{3 fases del ataque:}\label{fases-del-ataque}

{\def\LTcaptype{none} % do not increment counter
\begin{longtable}[]{@{}
  >{\raggedright\arraybackslash}p{(\linewidth - 4\tabcolsep) * \real{0.2500}}
  >{\raggedright\arraybackslash}p{(\linewidth - 4\tabcolsep) * \real{0.3333}}
  >{\raggedright\arraybackslash}p{(\linewidth - 4\tabcolsep) * \real{0.4167}}@{}}
\toprule\noalign{}
\begin{minipage}[b]{\linewidth}\raggedright
Fase
\end{minipage} & \begin{minipage}[b]{\linewidth}\raggedright
Acción
\end{minipage} & \begin{minipage}[b]{\linewidth}\raggedright
Detalles
\end{minipage} \\
\midrule\noalign{}
\endhead
\bottomrule\noalign{}
\endlastfoot
\textbf{1. Shutdown} & Apaga VMs &
\texttt{vim-cmd\ vmsvc/power.shutdown} \\
\textbf{2. Encrypt} & Encripta archivos & \texttt{.vmdk} (disco),
\texttt{.vmx} (config), \texttt{.vmsn} (snapshot) \\
\textbf{3. Exfil} & Roba datos & Para doble extorsión (encryption +
leak) \\
\end{longtable}
}

\subsection{Whitelist específico:}\label{whitelist-especuxedfico}

\begin{verbatim}
/bin, /etc, /usr, /lib, /var/log
\end{verbatim}

No encripta el sistema operativo base para que ESXi arranque.

\section{3. ALPHV Rust-based}\label{alphv-rust-based}

El libro analiza la \textbf{implementación Rust}:

\subsection{3.1. Configuración}\label{configuraciuxf3n}

Strings encriptados, \textbf{whitelisted directories} como \texttt{/bin}
y \texttt{/etc}.

\subsection{3.2. Proceso}\label{proceso}

\begin{enumerate}
\def\labelenumi{\arabic{enumi}.}
\tightlist
\item
  \textbf{Parar servicios}: \texttt{kill\ -9} procesos críticos
\item
  \textbf{Procesar archivos}: Encriptar con AES
\item
  \textbf{Eliminar shadow}: Borrar snapshots
\end{enumerate}

\subsection{3.3. Ventajas Rust}\label{ventajas-rust}

\begin{itemize}
\tightlist
\item
  \textbf{Memory safety} sin garbage collector
\item
  \textbf{Binarios estáticos} sin dependencias
\item
  \textbf{Performance} similar a C/C++
\item
  \textbf{Menos detección} por no usar APIs típicas de Windows
\end{itemize}

\section{3. ALPHV Rust-based}\label{alphv-rust-based-1}

Las nuevas variantes usan \textbf{Go/Rust}: - Binarios estáticos - Menos
detección heurística - Código nativo multiplataforma

\subsection{Características}\label{caracteruxedsticas}

\begin{Shaded}
\begin{Highlighting}[]
\CommentTok{// Whitelisting en Rust}
\KeywordTok{let}\NormalTok{ excluded }\OperatorTok{=} \PreprocessorTok{vec!}\NormalTok{[}\StringTok{"/bin"}\OperatorTok{,} \StringTok{"/etc"}\OperatorTok{,} \StringTok{"/usr"}\OperatorTok{,} \StringTok{"/lib"}\NormalTok{]}\OperatorTok{;}
\ControlFlowTok{if} \OperatorTok{!}\NormalTok{excluded}\OperatorTok{.}\NormalTok{iter()}\OperatorTok{.}\NormalTok{any(}\OperatorTok{|}\NormalTok{e}\OperatorTok{|}\NormalTok{ path}\OperatorTok{.}\NormalTok{starts\_with(e)) }\OperatorTok{\{}
\NormalTok{    encrypt\_file(path)}\OperatorTok{;}
\OperatorTok{\}}
\end{Highlighting}
\end{Shaded}

\section{4. Detección en Linux}\label{detecciuxf3n-en-linux}

\subsection{Syslog monitoring}\label{syslog-monitoring}

\begin{Shaded}
\begin{Highlighting}[]
\CommentTok{\# Auditd para monitorear encriptación masiva}
\ExtensionTok{{-}a}\NormalTok{ always,exit }\AttributeTok{{-}F}\NormalTok{ arch=b64 }\AttributeTok{{-}S}\NormalTok{ open,openat }\AttributeTok{{-}F}\NormalTok{ exit={-}EACCES }\AttributeTok{{-}k}\NormalTok{ suspicious\_file\_open}
\end{Highlighting}
\end{Shaded}

\subsection{Process monitoring}\label{process-monitoring}

\begin{Shaded}
\begin{Highlighting}[]
\CommentTok{\# Monitorear procesos sospechosos}
\FunctionTok{ps}\NormalTok{ aux }\KeywordTok{|} \FunctionTok{grep} \AttributeTok{{-}E} \StringTok{"(openssl|gpg|age)"}
\end{Highlighting}
\end{Shaded}

\section{5. Auto-Evaluación}\label{auto-evaluaciuxf3n-10}

\begin{enumerate}
\def\labelenumi{\arabic{enumi}.}
\tightlist
\item
  Diferencias clave Linux vs Windows ransomware
\item
  Técnicas de detección en ESXi
\item
  Ventajas del desarrollo en Rust
\end{enumerate}

\chapter{Laboratorio 11: Análisis de Ransomware
Linux/ESXi}\label{laboratorio-11-anuxe1lisis-de-ransomware-linuxesxi}

\chapter{Laboratorio 11: Análisis de Ransomware
Linux/ESXi}\label{laboratorio-11-anuxe1lisis-de-ransomware-linuxesxi-1}

\section{Objetivo}\label{objetivo-10}

Analizar variantes de ransomware para sistemas Unix sin ejecutar código.

\section{Paso 1: Identificar paths de
ESXi}\label{paso-1-identificar-paths-de-esxi}

\subsection{1.1. Configuración típica}\label{configuraciuxf3n-tuxedpica}

\begin{Shaded}
\begin{Highlighting}[]
\FunctionTok{strings}\NormalTok{ esxi\_ransomware.elf }\KeywordTok{|} \FunctionTok{grep} \AttributeTok{{-}i} \AttributeTok{{-}E} \StringTok{"(vmfs|vmdk|vmx|vsphere)"}
\end{Highlighting}
\end{Shaded}

\section{Paso 2: Identificar rutas
excluidas}\label{paso-2-identificar-rutas-excluidas}

\subsection{2.1. Whitelist Linux}\label{whitelist-linux}

\begin{Shaded}
\begin{Highlighting}[]
\FunctionTok{strings}\NormalTok{ ransomware.elf }\KeywordTok{|} \FunctionTok{grep} \AttributeTok{{-}i} \AttributeTok{{-}E} \StringTok{"(bin|etc|usr|lib|systemd)"}
\end{Highlighting}
\end{Shaded}

\section{Paso 3: Identificar procesos
atacados}\label{paso-3-identificar-procesos-atacados}

\subsection{3.1. Process killing}\label{process-killing}

\begin{Shaded}
\begin{Highlighting}[]
\FunctionTok{strings}\NormalTok{ ransomware.elf }\KeywordTok{|} \FunctionTok{grep} \AttributeTok{{-}i} \AttributeTok{{-}E} \StringTok{"(kill|shutdown|pkill|systemctl)"}
\end{Highlighting}
\end{Shaded}

\section{Verificación}\label{verificaciuxf3n-9}

{\def\LTcaptype{none} % do not increment counter
\begin{longtable}[]{@{}ll@{}}
\toprule\noalign{}
Checkpoint & Resultado \\
\midrule\noalign{}
\endhead
\bottomrule\noalign{}
\endlastfoot
VMFS paths & Identificado \\
Whitelist & Detectado \\
Process kill & Encontrado \\
\end{longtable}
}

\chapter{Quiz 11: Linux y ESXi
Ransomware}\label{quiz-11-linux-y-esxi-ransomware}

\chapter{Quiz 11: Linux y ESXi
Ransomware}\label{quiz-11-linux-y-esxi-ransomware-1}

\subsection{1. Identificar}\label{identificar-6}

Extensión común de archivo de disco de VM.

\subsection{2. Explicar}\label{explicar-6}

Ventaja de Rust para malware.

\subsection{3. Analizar}\label{analizar-3}

Técnica de REvil para VM shutdown.

\subsection{4. Verificar}\label{verificar-4}

Método para monitorear encriptación en Linux.

\subsection{5. Comando}\label{comando-9}

Escribe el command para listar procesos sospechosos.

\part{Unidad 12: Lecciones de la Economía Underground}

\chapter{Lecciones de la Economía
Underground}\label{lecciones-de-la-economuxeda-underground}

\chapter{Unidad 12: Lecciones de la Economía
Underground}\label{unidad-12-lecciones-de-la-economuxeda-underground-1}

\section{1. Introducción}\label{introducciuxf3n-11}

La dark web no es caos: es un \textbf{ecosistema económico} con
dinámicas claras. Entender estas dinámicas ayuda a predecir la próxima
ola de amenazas.

\section{2. Modelo económico del
underground}\label{modelo-econuxf3mico-del-underground}

\subsection{2.1. Reactive Development}\label{reactive-development}

Los actores reaccionan a: - Parches de seguridad - Operaciones de law
enforcement - Nuevas herramientas de defensa

\subsection{2.2. Rebranding}\label{rebranding}

Productos criminales: - Cambio de nombre (REvil → Sodinokibi) -
Actualización de versiones - Nuevos ``vendedores''

\section{3. Lowered Barriers to Entry}\label{lowered-barriers-to-entry}

\subsection{3.1. DIY malware kits}\label{diy-malware-kits}

Cualquiera puede comprar: - Phishing kits (\$50-\$500) - Loaders
(\$20-\$200) - Ransomware RaaS (30\% profit share)

\subsection{3.2. Tutoriales publicados}\label{tutoriales-publicados}

La dark web es \textbf{educativa}: - Tutoriales de pentesting - Manuales
de uso de herramientas - ``Cursos'' de ciberdelincuencia

\section{4. Inteligencia Artificial}\label{inteligencia-artificial}

Aplicaciones emergentes: - \textbf{Generación de phishing emails} con
LLMs - \textbf{Evasión de antivirus} con prompts - \textbf{Análisis
automático} de vulnerabilidades

\section{5. Futuro de la
ciberseguridad}\label{futuro-de-la-ciberseguridad}

\subsection{Tendencias}\label{tendencias}

\begin{enumerate}
\def\labelenumi{\arabic{enumi}.}
\tightlist
\item
  \textbf{IA generativa} como arma y defensa
\item
  \textbf{Rust/Go} para malware más eficiente
\item
  \textbf{Ransomware} contra infra crítica
\item
  \textbf{Deepfakes} para ingeniería social
\end{enumerate}

\subsection{Estrategias defensivas}\label{estrategias-defensivas}

\begin{itemize}
\tightlist
\item
  Threat hunting proactivo
\item
  Zero-trust architecture
\item
  AI-powered detection
\item
  Supply chain security
\end{itemize}

\section{6. Proyecto Final}\label{proyecto-final}

Analiza un caso real del underground (sin acceder) y escribe: -
Herramientas involucradas - Kill chain del ataque - Mitigaciones
propuestas - Impacto económico

\section{7. Auto-Evaluación}\label{auto-evaluaciuxf3n-11}

\begin{enumerate}
\def\labelenumi{\arabic{enumi}.}
\tightlist
\item
  ¿Qué significa reactive development?
\item
  Modelo de negocio de Ransomware-as-a-Service
\item
  Impacto de IA en el underground
\end{enumerate}

\chapter{Laboratorio 12: Proyecto
Final}\label{laboratorio-12-proyecto-final}

\chapter{Laboratorio 12: Proyecto
Final}\label{laboratorio-12-proyecto-final-1}

\section{Objetivo}\label{objetivo-11}

Completar un análisis integral de un caso real del underground y
presentar recomendaciones de defensa.

\section{Paso 1: Selección de caso}\label{paso-1-selecciuxf3n-de-caso}

\subsection{1.1. Fuentes permitidas}\label{fuentes-permitidas}

\begin{itemize}
\tightlist
\item
  Reportes de Symantec/Mandiant
\item
  GitHub: DissectingTheDarkWeb
\item
  Blogposts de investigadores públicos
\end{itemize}

\textbf{NO} acceder a foros reales

\section{Paso 2: Análisis del caso}\label{paso-2-anuxe1lisis-del-caso}

\subsection{2.1. Kill chain mapping}\label{kill-chain-mapping}

{\def\LTcaptype{none} % do not increment counter
\begin{longtable}[]{@{}lll@{}}
\toprule\noalign{}
Fase Kill Chain & Técnica observada & Herramienta usada \\
\midrule\noalign{}
\endhead
\bottomrule\noalign{}
\endlastfoot
1. Reconocimiento & & \\
2. Arma & & \\
3. Entrega & & \\
4. Explotación & & \\
5. Instalación & & \\
6. C2 & & \\
7. Acción & & \\
\end{longtable}
}

\section{Paso 3: Recomendaciones}\label{paso-3-recomendaciones}

\subsection{3.1. Mitigación por fase}\label{mitigaciuxf3n-por-fase}

Para cada fase identifica: - Señales de detección - Controles
preventivos - Respuesta a incidentes

\section{Entregable}\label{entregable}

\begin{enumerate}
\def\labelenumi{\arabic{enumi}.}
\tightlist
\item
  Reporte de 2-3 páginas (syllabus + análisis + mitigaciones)
\item
  Presentación de 5 diapositivas
\item
  Lista de IoCs recopilados
\end{enumerate}

\section{Rubrica}\label{rubrica-1}

{\def\LTcaptype{none} % do not increment counter
\begin{longtable}[]{@{}llll@{}}
\toprule\noalign{}
Criterio & Excelente & Bueno & Mejorable \\
\midrule\noalign{}
\endhead
\bottomrule\noalign{}
\endlastfoot
Análisis & Kill chain completo & 5/7 fases & \textless5 fases \\
IoCs & 10+ identificados & 5-10 & \textless5 \\
Mitigaciones & Al menos 3 & 1-2 & 0-1 \\
Presentación & Clara y técnica & Básica & Confusa \\
\end{longtable}
}

\chapter{Quiz 12: Lecciones de la Economía
Underground}\label{quiz-12-lecciones-de-la-economuxeda-underground}

\chapter{Quiz 12: Lecciones de la Economía
Underground}\label{quiz-12-lecciones-de-la-economuxeda-underground-1}

\subsection{1. Definir}\label{definir-1}

¿Qué significa reactive development en threat actors?

\subsection{2. Explicar}\label{explicar-7}

Modelo de negocio de Ransomware-as-a-Service.

\subsection{3. Analizar}\label{analizar-4}

Impacto de IA generativa en el underground.

\subsection{4. Proponer}\label{proponer}

Una mitigación para futuras amenazas de stealers.

\subsection{5. Reflexionar}\label{reflexionar}

¿Cómo cambia el threat landscape con Rust malware?




\end{document}
